\documentclass[11pt]{article}
\usepackage[margin=1in]{geometry}
\usepackage[T1]{fontenc}
\usepackage{amsmath,amssymb,amsthm,mathtools}
\usepackage{booktabs,array}
\usepackage{xcolor}
\usepackage{fancyvrb}
\usepackage[hidelinks]{hyperref}
\usepackage[numbers,sort&compress]{natbib}

\newcommand{\Z}{\mathbb{Z}}
\newcommand{\C}{\mathbb{C}}
\newcommand{\E}{\mathbb{E}}
\newcommand{\tr}{\operatorname{tr}}
\newcommand{\mult}{\operatorname{mult}}
\newcommand{\Sym}{\operatorname{Sym}}
\newcommand{\Ptor}{P_{\mathrm{tor}}}
\newcommand{\ket}[1]{\lvert #1\rangle}
\newcommand{\bra}[1]{\langle #1\rvert}

\newtheorem{theorem}{Theorem}
\newtheorem{lemma}[theorem]{Lemma}
\newtheorem{proposition}[theorem]{Proposition}
\newtheorem{corollary}[theorem]{Corollary}
\newtheorem{conjecture}[theorem]{Conjecture}
\theoremstyle{definition}
\newtheorem{definition}[theorem]{Definition}
\theoremstyle{remark}
\newtheorem{remark}[theorem]{Remark}

\title{Sidon Sets, Additive Energy, and the Optimal Cloning\\ of Finite Phase-State Codebooks}
\author{Michael Weldon\\[2pt]
\small Independent Researcher, Las Vegas, NV, USA\\
\small ORCID: \href{https://orcid.org/0009-0001-6119-0560}{0009-0001-6119-0560}}
\date{Version 2 --- 18 October 2026}

\begin{document}
\maketitle

\begin{abstract}
Phase states $\ket{\psi_x}=d^{-1/2}\sum_{j}\omega^{x k_j}\ket{j}$, indexed by $x\in\Z_q$ and a key set
$K=\{k_1,\dots,k_d\}\subset\Z_q$, are the state family behind quantum fingerprinting and quantum hashing.
We ask how well an adversary holding $T$ copies of $\ket{\psi_x}$ can produce $T+k$ copies, and when a
\emph{finite} codebook is exactly as hard to clone as the continuous family of all phase states (the torus),
whose optimal cloning fidelity was recently determined. The answer is governed by additive combinatorics of $K$.
(i) The optimal $1\to2$ global cloning fidelity equals $(2d-1)/d^2$ if and only if $K$ is a Sidon set.
(ii) A $B_{T+k}$ key set attains the torus optimum: analytically for $d=2$ (all $T$) and for $T=1$ (all $d,k$), and for
$T=2,3$ (all $d,k$) by a proof whose finitely many computational steps are verified by standalone exact and
interval-arithmetic checkers; we conjecture it for all $T$. (iii) For $T=1$ we characterise exactly, for every $k$, the key
sets that attain the torus optimum; the $B_{1+k}$ condition is sufficient but not necessary. (iv) The optimal fidelity of
fusing $m$ phase states into the state of their sum equals $E_m(K)/d^{m+1}$, where $E_m$ is the higher additive energy;
Sidon sets minimise it for every $m$. Small Fourier bias, the standard design criterion for quantum hashing, does not
control cloning. For finite orbits of non-abelian groups we prove a window theorem giving a dual-feasibility criterion, an
obstruction, and a representation-theoretic selection rule for when matching the $(T+k)$-th moment suffices, and we give
certified counterexamples for $SU(3)$ and $SU(2)$ showing that moment matching alone does not.
\end{abstract}
\noindent\textbf{Keywords:} quantum cloning, phase states, quantum hashing, Sidon sets, $B_h$ sets, additive energy.\\
\textbf{MSC2020:} 81P45 (primary); 11B13, 05B10 (secondary).

\section{Introduction}
Quantum fingerprints~\cite{Buhrman2001} and quantum hash functions~\cite{AblayevVasiliev2014,Vasiliev2016,Ziiatdinov2016,Ablayev2025}
map a classical string $x$ to a short quantum state. In the phase form, $x\in\Z_q$ is encoded as
\begin{equation}\label{eq:phase}
\ket{\psi_x}=\frac{1}{\sqrt d}\sum_{j=1}^{d}\omega^{x k_j}\ket{j},\qquad \omega=e^{2\pi i/q},\qquad K=\{k_1,\dots,k_d\}\subset\Z_q .
\end{equation}
The standard design criterion is that $K$ be $\varepsilon$-biased, which bounds the overlaps
$|\langle\psi_x|\psi_y\rangle|$ and hence collision resistance~\cite{Vasiliev2016,Ablayev2025}.
When such states serve as quantum public keys or tokens~\cite{GottesmanChuang2001}, a second property matters:
how well can an adversary holding $T$ copies \emph{counterfeit} $T+k$ copies?

For the continuous family of all phase states, $\ket{\psi_\theta}=d^{-1/2}\sum_j e^{i\theta_j}\ket{j}$ with
$\theta$ uniform on the torus, the optimal $T\to T+k$ cloning fidelity is known: special cases go back to
phase-covariant cloning~\cite{Buscemi2005,KimChitambar2022}, and the general value was recently obtained by
Herasymenko, Lin, Verduyn Lunel and Grinko~\cite{Herasymenko2026}. A finite codebook is a subset of the torus, so it can
only be easier to clone. The question we answer is \emph{when it is exactly as hard}.

\paragraph{Results.}
\begin{itemize}
\item \textbf{Reduced dual (Lemma~\ref{lem:R}).} For every $K$ and $(T,k)$ the counterfeiting SDP reduces to a
small convex program in the representation counts of $K$.
\item \textbf{Sidon theorem (Theorem~\ref{thm:sidon}).} $P_K(1\to2)=(2d-1)/d^2$ if and only if $K$ is Sidon.
\item \textbf{Torus theorem (Theorem~\ref{thm:torus}).} For a $B_{T+k}$ key set, $P_K(T\to T+k)=\Ptor$ in the proved cases
$d=2$ (all $T$) and $T\le3$ (all $d,k$); conjectured in general (Conjecture~\ref{conj:TT}), with exact certificates on the
finite boxes listed after the conjecture.
\item \textbf{Fusion theorem (Theorem~\ref{thm:fusion}).} The optimal fusion fidelity is the additive energy $E_q(K)/d^3$.
\item \textbf{Exact characterisation at $T=1$ (Theorem~\ref{thm:t1iff}).} For every $k$, $P_K(1,k)=\Ptor$ if and only if no additive class
of $K$ carries more weight than the torus maximum. $B_{1+k}$ is sufficient but not necessary: Sidon sets that are not $B_3$ can attain the torus
value (Section~\ref{sec:t1iff}).
\item \textbf{Multi-fold fusion (Theorem~\ref{thm:fusionm}).} The optimal fidelity of fusing $m$ phase states is the higher energy $E_m(K)/d^{m+1}$;
Sidon sets minimise it for every $m$.
\item \textbf{Negative results.} Small bias does not imply counterfeit resistance; the Sidon sets in our tables that are not $B_3$ fail at $T+k=3$,
and arithmetic progressions are worst in most cases tested (Section~\ref{sec:neg}).
\item \textbf{Non-abelian orbits (Section~\ref{sec:nonab}).} A window theorem for finite orbits $X=G\cdot f\subset H\cdot f$,
and certified counterexamples ($SU(3)$ on $\Sym^2\C^3$; $SU(2)$ spin 1) to the naive moment-matching principle.
\end{itemize}
The torus values themselves are not new: they are given in full by~\cite{Herasymenko2026}, with special cases
in~\cite{Buscemi2005,KimChitambar2022}. What is new is the exact characterisation of the \emph{finite} codebooks that attain them, the
additive-energy form of optimal fusion, and the failure of the corresponding moment-matching principle for finite orbits of
non-abelian groups. Section~\ref{sec:related} discusses prior work in detail.

\section{Setting and notation}\label{sec:setting}
Throughout, $K=(k_1,\dots,k_d)$ consists of $d\ge2$ distinct elements of $\Z_q$, and $x\in\Z_q$ is uniform.

\begin{definition}[Counterfeiting value]
For $T\ge1$, $k\ge1$,
\[
P_K(T,k)=\max_{\Phi}\ \E_{x\in\Z_q}\,\bra{\psi_x^{\otimes(T+k)}}\,\Phi\big(\psi_x^{\otimes T}\big)\,\ket{\psi_x^{\otimes(T+k)}},
\]
the maximum over quantum channels (CPTP maps) $\Phi$ from $\Sym^T(\C^d)$ to $\Sym^{T+k}(\C^d)$; this is the
average \emph{global} fidelity with the ideal output $\psi_x^{\otimes(T+k)}$. The torus value $\Ptor(d;T,k)$ is the same
quantity for $\ket{\psi_\theta}=d^{-1/2}\sum_je^{i\theta_j}\ket j$ with $\theta$ uniform on $[0,2\pi)^d$. Twirling any channel over
the phase group $\theta\mapsto\theta+\phi$ preserves the average fidelity and produces a covariant channel, whose fidelity is the
same on every phase state; hence the optimal average equals the optimal worst-case fidelity $F^\star_{\rm ph}(d,T,k)$
of~\cite[Thm.~II.3; formal statement Thm.~X.1, Eq.~(16)]{Herasymenko2026}.
\end{definition}

\begin{definition}[$B_h$ sets, representation counts, additive energy]
$K$ is a $B_h$ set (mod $q$) if any two distinct \emph{multisets} of size $h$ drawn from $K$ (repetition allowed, order
ignored) have distinct sums mod $q$; $B_2$ sets are \emph{Sidon} sets~\cite{ErdosTuran1941,OBryant2004}. A $B_h$ set is $B_{h'}$
for every $h'\le h$ (add the same $h-h'$ elements to both multisets). By contrast, the counts $r_m$ below are over \emph{ordered}
tuples. The ordered representation count is
$r_m(v)=\#\{(j_1,\dots,j_m)\in[d]^m:\sum_i k_{j_i}\equiv v\}$, and the additive energy is
$E_q(K)=\#\{(a,b,c,e)\in K^4:a+e\equiv b+c \pmod q\}$.
\end{definition}

We write $D_m=\binom{d+m-1}{m}=\dim\Sym^m(\C^d)$ and $W(d,T,k)=D_T/D_{T+k}$, Werner's optimal universal cloning
fidelity~\cite{Werner1998}. For multi-indices $\gamma\in\Z_{\ge0}^d$, $\mult(\gamma)=|\gamma|!/\prod_i\gamma_i!$.

\section{The reduced dual}\label{sec:dual}
The counterfeiting SDP has dual $P_K(T,k)=\min\{\tr Y : Y\otimes I\succeq X\}$, with
$X=\E_x \bar\psi_x^{\otimes T}\otimes\psi_x^{\otimes (T+k)}$ restricted to $\Sym^T\otimes\Sym^{T+k}$.

\begin{lemma}[Reduced dual]\label{lem:R}
For every $K$ and $(T,k)$,
\[
P_K(T,k)=\min_{z\in(0,\infty)^{\mathrm{supp}\,r_T}}\ \sum_{w\in\mathrm{supp}\,r_T}\frac{r_T(w)}{z_w}\quad\text{s.t.}\quad d^{-(2T+k)}\sum_{w} r_{T+k}(s+w)\,z_w\le1\qquad\forall s\in\Z_q ,
\]
the sums running over $w\in\mathrm{supp}\,r_T$; the minimum is attained.
\end{lemma}
\begin{proof}[Proof sketch; full proof in Appendix~\ref{app:lemR}]
In occupation bases $\ket\alpha$ ($|\alpha|=T$) and $\ket\beta$ ($|\beta|=T+k$), with $c_\gamma=\sqrt{\mult(\gamma)}$,
averaging over $x$ gives $X=\sum_{s\in\Z_q}\ket{w_s}\bra{w_s}$ with
$w_s=d^{-(2T+k)/2}\sum_{K\cdot(\beta-\alpha)\equiv s}c_\alpha c_\beta\ket{\alpha,\beta}$. $X$ and the constraint are
invariant under $\bar U_x^{\otimes T}\otimes U_x^{\otimes(T+k)}$, $U_x=\mathrm{diag}(\omega^{xk_j})$, so an optimal $Y$
may be taken block diagonal over the classes $\{\alpha:K\cdot\alpha\equiv w\}$. Within sector $s$, $Y\otimes I\succeq X$
holds iff $\sum_{(\beta,w)}c_\beta^2 d^{-(2T+k)}v_w^\dagger Y_w^{-1}v_w\le1$ with $v_w(\alpha)=c_\alpha$. For fixed
$g_w=v_w^\dagger Y_w^{-1}v_w$, Cauchy--Schwarz gives $\tr Y_w\ge\lVert v_w\rVert^2/g_w=r_T(w)/g_w$, attained by a rank-one
$Y_w$; put $z_w=g_w$ and use $\sum_{K\cdot\beta\equiv v}c_\beta^2=r_{T+k}(v)$. Strong duality holds by Slater.
\end{proof}

\begin{lemma}\label{lem:lower}
For every $K$, $P_K(T,k)\ge\Ptor(d;T,k)$.
\end{lemma}
\begin{proof}
The optimal torus cloner of~\cite{Herasymenko2026} is covariant and achieves fidelity $\Ptor$ on every phase state, in
particular on every $\ket{\psi_x}$.
\end{proof}

\begin{lemma}[Moment bound]\label{lem:B}
For any finite ensemble $\{\psi_x\}$, $P\le \dim\big(\mathrm{span}\{\psi_x^{\otimes T}\}\big)\,\lambda_{\max}\big(\E_x\psi_x^{\otimes(T+k)}\big)$.
\end{lemma}
\begin{proof}
$\bar\psi^{\otimes T}\preceq\Pi_V$ on the span $V$ of the inputs, so $X\preceq\Pi_V\otimes M$ with
$M=\E\psi^{\otimes(T+k)}$, and $\tr(JX)\le\lambda_{\max}(M)\tr\big((\Pi_V\otimes I)J\big)=\lambda_{\max}(M)\dim V$ for every Choi
operator $J$ with $\tr_{\rm out}J=I$.
\end{proof}

\begin{proposition}[Near-optimality of $B_{T+k}$ codebooks, all $T$]\label{prop:nearW}
If $K$ is $B_{T+k}$ and $d\ge T+k$, then
$\Ptor\le P_K(T,k)\le W(d,T,k)\prod_{i=0}^{T+k-1}\big(1+\tfrac id\big)$.
\end{proposition}
\begin{proof}
Under $B_{T+k}$, $\E_x\psi_x^{\otimes(T+k)}$ is diagonal in the occupation basis with entries $\mult(\beta)/d^{T+k}$, whose
maximum is $(T+k)!/d^{T+k}$ for $d\ge T+k$. Apply Lemma~\ref{lem:B} with $\dim V\le D_T$ and use $D_m=\prod_{i<m}(d+i)/m!$.
\end{proof}

\section{The Sidon theorem}\label{sec:sidon}
\begin{theorem}\label{thm:sidon}
For every $K\subset\Z_q$ with $|K|=d\ge2$, $P_K(1,1)=\dfrac{2d-1}{d^2}$ if and only if $K$ is a Sidon set; otherwise
$P_K(1,1)>(2d-1)/d^2$.
\end{theorem}
\begin{proof}
\emph{Primal.} The channel with Kraus operators $A_j\ket a=\sqrt{(2-\delta_{aj})/(2d-1)}\,\ket{e_a+e_j}_{\Sym^2}$, $j\in[d]$, is
trace preserving and has fidelity $(2d-1)/d^2$ on every phase state, so $P_K(1,1)\ge(2d-1)/d^2$ for all $K$.

\emph{Dual.} In Lemma~\ref{lem:R} with $T=k=1$, take $z_a\equiv d^3/(2d-1)$ (i.e.\ $Y=\frac{2d-1}{d^3}I$, trace $(2d-1)/d^2$). Torus
sectors are indexed by charges $c=\beta-\alpha$. Type-1 charges $c=e_j$ have weight $\sum_a d^{-3}(2-\delta_{aj})=(2d-1)/d^3$ and
are tight; type-2 charges $c=e_b+e_{b'}-e_a$ ($a\notin\{b,b'\}$) have value $(2-\delta_{bb'})/(2d-1)$. For a Sidon set, a type-1
sector merges with nothing (two type-1 charges merge only if $k_i\equiv k_j$; type 1 with type 2 only through a non-trivial
relation $k_j+k_a\equiv k_b+k_{b'}$). In a merged class of type-2 charges, Sidon makes the representation $s+k_a\equiv k_b+k_{b'}$
unique for each $a$, so the class value is at most $\sum_a(2-\delta_{b_ab'_a})/(2d-1)$, which can exceed 1 only if all $d$ edges
$E_a=\{b_a,b'_a\}$ avoid $a$. If two such edges share a vertex, $E_a=\{u,v\}$, $E_{a'}=\{u,w\}$, subtracting gives
$k_a+k_w\equiv k_{a'}+k_v$, and Sidon forces $a=v\in E_a$, a contradiction; hence the $d$ edges are pairwise disjoint on
$d$ vertices, which is impossible. So $Y$ is feasible.

\emph{Only if.} We use the Lagrangian lower bound for the convex program of Lemma~\ref{lem:R}. With $T=k=1$ and distinct keys,
$r_1(w)=1$ for $w\in K$, so $P_K(1,1)=\min_{z>0}\sum_a 1/z_a$ subject to $g_s(z):=d^{-3}\sum_a r_2(s+k_a)z_a\le1$ for all $s$. For any
$\lambda\ge0$ supported on $\{s\in K\}$ and any feasible $z$,
\[
\sum_a\frac1{z_a}\ \ge\ \sum_a\frac1{z_a}+\sum_{s}\lambda_s\big(g_s(z)-1\big)\ \ge\ \min_{z>0}\Big[\sum_a\Big(\frac1{z_a}+c_a z_a\Big)\Big]-\sum_s\lambda_s
=2\sum_a\sqrt{c_a}-\sum_s\lambda_s ,
\]
with $c_a=d^{-3}\sum_{s}\lambda_s r_2(s+k_a)$. Take $\lambda_s=\lambda:=(2d-1)/d^3$ for every $s\in K$. For every set of distinct keys,
$r_2(k_j+k_a)\ge2-\delta_{ja}$, so $c_a\ge\lambda(2d-1)/d^3=\lambda^2$, with equality for all $a$ exactly when $K$ is Sidon (then the
bound equals $2d\lambda-d\lambda=(2d-1)/d^2$). If $K$ is not Sidon there are multisets $\{i,j\}\ne\{l,m\}$ with
$k_i+k_j\equiv k_l+k_m$; then $r_2(k_i+k_j)>2-\delta_{ij}$, hence $c_i>\lambda^2$ and $P_K(1,1)\ge2\sum_a\sqrt{c_a}-d\lambda>(2d-1)/d^2$.
\end{proof}

\begin{remark}
The value $(2d-1)/d^2$ itself is the $1\to2$ torus optimum, known from~\cite{KimChitambar2022} (and the $N=1$ case
of~\cite[Eq.~(16)]{Herasymenko2026}); the content of Theorem~\ref{thm:sidon} is the exact characterisation of the finite codebooks that
attain it. For $d=2$ it recovers the observation of~\cite{Jing2012} that three equally spaced equatorial qubit states already
determine the phase-covariant cloner (there for single-copy fidelity). Our necessity statements concern \emph{global} fidelity only.
\end{remark}

\section{The torus theorem}\label{sec:torus}
\begin{theorem}[Torus theorem]\label{thm:torus}
Let $K$ be a $B_{T+k}$ set mod $q$. Then $P_K(T,k)=\Ptor(d;T,k)$ in each of the following cases:
(a) $d=2$, all $T,k\ge1$;\quad (b) $T\in\{1,2,3\}$, all $d\ge2$, $k\ge1$.
\end{theorem}

\subsection{Reduction to a packing inequality}
Let $C$ be the set of balanced vectors with entries $\lfloor k/d\rfloor$ or $\lfloor k/d\rfloor+1$ summing to $k$, and for
$|b|=T$ put $Z_b=\sum_{g\in C}\mult(b+g)$ and $a_b=\sqrt{\mult(b)/Z_b}$. For a shift $\Delta\in\Z^d$ with $\sum_i\Delta_i=k$ define
\[
\Lambda(\Delta)=\sum_{|b|=T,\ b\ge\Delta^-}a_b\,\mult(b+\Delta),\qquad \Lambda_0=\Lambda(g)\ (g\in C),\qquad
\rho(F)=\frac{1}{\Lambda_0}\sum_{\Delta\in F}\Lambda(\Delta).
\]
\begin{lemma}[Separation]\label{lem:sep}
Let $K$ be a $B_h$ set mod $q$ and let $w\in\Z^d$, $w\ne0$, satisfy $\sum_iw_i=0$ and $\sum_ik_iw_i\equiv0\pmod q$. Then
$\lVert w\rVert_1\ge2(h+1)$.
\end{lemma}
\begin{proof}
Since $\sum_iw_i=0$, the positive and negative parts $w^\pm\ge0$ satisfy $|w^+|=|w^-|=\lVert w\rVert_1/2$, and they have disjoint
supports, so they are distinct multisets on $[d]$ ($w\ne0$). If $\lVert w\rVert_1\le2h$, add the same multiset of size $h-|w^+|$
(e.g.\ copies of index $1$) to both; this gives two distinct multisets of size $h$ whose $K$-sums differ by $\sum_ik_iw_i\equiv0$,
contradicting $B_h$.
\end{proof}

\begin{proposition}\label{prop:reduction}
If $K$ is $B_{T+k}$, two distinct torus charges merged into the same class of Lemma~\ref{lem:R} satisfy
$\lVert\Delta-\Delta'\rVert_1\ge2(T+k+1)$. Consequently, if every family $F$ of shifts $\Delta\in\Z^d$ with $\sum_i\Delta_i=k$ and $|\Delta^-|\le T$ (the charges $\beta-\alpha$ with
$|\alpha|=T$, $|\beta|=T+k$), pairwise $\lVert\Delta-\Delta'\rVert_1\ge 2(T+k+1)$ and $|F|\ge2$ satisfies $\rho(F)\le1$, then the torus dual point
$z^\star_\alpha=d^{2T+k}a_\alpha/\Lambda_0$ (which corresponds to the dual point of~\cite[Eq.~(147)]{Herasymenko2026} via $y_\alpha=\mult(\alpha)/z_\alpha$) is feasible for $K$, and $P_K(T,k)=\Ptor$.
\end{proposition}
\begin{proof}
\emph{Classes.} $B_{T+k}$ implies $B_T$, so each class $\{\alpha:K\cdot\alpha\equiv w\}$ of Lemma~\ref{lem:R} is a single multiset
$\alpha$, with $r_T(w)=\mult(\alpha)$; likewise $r_{T+k}(v)=\mult(\beta)$ for the unique $\beta$ with $K\cdot\beta\equiv v$, if any. The
constraint of class $s$ therefore reads $\sum_{(\alpha,\beta):K\cdot(\beta-\alpha)\equiv s}\mult(\beta)\,z_\alpha\le d^{2T+k}$. Group the
pairs by their charge $\Delta=\beta-\alpha$; the charges occurring form the family $F_s=\{\Delta:K\cdot\Delta\equiv s\}$, and every $\alpha$
with $\alpha+\Delta\ge0$ occurs with each $\Delta\in F_s$. Two distinct $\Delta,\Delta'\in F_s$ have $w=\Delta-\Delta'\ne0$, $\sum w_i=0$ and
$K\cdot w\equiv0$, so $\lVert\Delta-\Delta'\rVert_1\ge2(T+k+1)$ by Lemma~\ref{lem:sep}. The class-$s$ constraint is
$\sum_{\Delta\in F_s}\sum_{\alpha\ge\Delta^-}\mult(\alpha+\Delta)\,z_\alpha\le d^{2T+k}$.

\emph{Torus dual.} Put $z^\star_\alpha=d^{2T+k}a_\alpha/\Lambda_0$. Then the class-$s$ constraint at $z^\star$ is exactly $\rho(F_s)\le1$. For
$|F_s|=1$ this is $\Lambda(\Delta)\le\Lambda_0$, which is the feasibility of $z^\star$ in the torus problem (where every class is a single
charge); for $|F_s|\ge2$ it is the hypothesis. The objective is
$\sum_\alpha\mult(\alpha)/z^\star_\alpha=d^{-(2T+k)}\Lambda_0\sum_\alpha\sqrt{\mult(\alpha)Z_\alpha}$. Since $a$ is $S_d$-invariant and $C$ is a single
$S_d$-orbit, $\Lambda(g)=\Lambda_0$ for all $g\in C$, and averaging over $C$ gives $\Lambda_0=|C|^{-1}\sum_\alpha a_\alpha Z_\alpha
=|C|^{-1}\sum_\alpha\sqrt{\mult(\alpha)Z_\alpha}$. Hence the objective equals $|C|\Lambda_0^2/d^{2T+k}$, which is exactly the optimal fidelity
$F^\star_{\rm ph}(d,T,k)=\Ptor$ of~\cite[Eq.~(16)]{Herasymenko2026}. Therefore $P_K(T,k)\le\Ptor$, and Lemma~\ref{lem:lower} gives equality.
\end{proof}
\begin{remark}[Identification with the dual certificate of~\cite{Herasymenko2026}]\label{rem:torusdual}
In the notation of~\cite[Sec.~X]{Herasymenko2026} ($N=T$, $K=k$, $\mathcal C_{d,K}=C$, $\mathcal D_\Delta=\{b:b+\Delta\ge0\}$ by their Eq.~(111)), the
phase-covariant dual of their Prop.~X.2 has variables $y_b>0$ and, by their Eq.~(119), one scalar constraint per shift,
$\sum_{b\in\mathcal D_\Delta}|v_\Delta(b)|^2/y_b\le1$ with $|v_\Delta(b)|^2=d^{-(2T+k)}\mult(b)\mult(b+\Delta)$ (their Eq.~(115)). Their quantities $Z_b$
(Eq.~(101)) and $\Lambda(\Delta)$ (Eq.~(141)) coincide with ours. Under $y_b=\mult(b)/z_b$ this constraint is exactly the single-charge
constraint $d^{-(2T+k)}\sum_{b\ge\Delta^-}\mult(b+\Delta)z_b\le1$ of Lemma~\ref{lem:R}, and the objectives agree, $\sum_by_b=\sum_b\mult(b)/z_b$. Their dual
point, Eq.~(147), $y_b=\Lambda(\Gamma_0)\sqrt{\mult(b)Z_b}/d^{2T+k}$, corresponds to $z_b=\mult(b)/y_b=d^{2T+k}a_b/\Lambda_0=z^\star_b$. Their Eq.~(148)
shows that the constraint of shift $\Delta$ at this point equals $\Lambda(\Delta)/\Lambda(\Gamma_0)$, and their Prop.~X.5 (Eq.~(146)) gives
$\Lambda(\Delta)\le\Lambda(\Gamma_0)=\Lambda_0$ for every integer shift with $|\Delta|=k$. Hence $z^\star$ is feasible for every single-charge class,
which is the step used in the proof above; its objective is their Eq.~(150), $|C|\Lambda_0^2/d^{2T+k}=F^\star_{\rm ph}(d,T,k)$.
\end{remark}

\subsection{The packing inequality}
\begin{theorem}\label{thm:packing}
Every family as in Proposition~\ref{prop:reduction} with $|F|\ge2$ satisfies:
(a) $\rho(F)\le1$ for $d=2$, all $T,k$;
(b) $\rho(F)\le\frac{d}{(k+1)(d-k/2)}$ for $T=1$, $k\le d-1$, and $\rho(F)\le\frac{3}{4(d-1)}$ for $T=1$, $k\ge d$;
(c) $\rho(F)\le0.9961$ for $T=2$, $d\ge3$;
(d) $\rho(F)\le0.9706$ for $T=3$, $d\ge3$.
\end{theorem}
Parts (a) and (b) have complete pen-and-paper proofs (Appendix~\ref{app:proofs}). Parts (c) and (d) are computer-assisted:
analytic bounds (normalisation Lemma N, structure Lemma S, multinomial moves Lemma M, and a Johnson-type multiset packing
bound Lemma J) reduce them to finitely many closed-form evaluations, which are verified in interval arithmetic, and to finitely
many residual cases closed separately (13 at $T=2$, by Lagrangian bounds evaluated in interval arithmetic; 18 at $T=3$, by exact
clique-LP certificates); see Appendix~\ref{app:ca}. The decimal constants are
outward-rounded upper bounds; every certificate is compared with $1$ in exact rational or interval arithmetic.

\begin{conjecture}\label{conj:TT}
For all $d\ge3$, $T\ge4$ and $k\ge1$, every family as in Proposition~\ref{prop:reduction} with $|F|\ge2$ satisfies $\rho(F)<1$;
in particular, $P_K(T,k)=\Ptor$ for every $B_{T+k}$ set $K$.
\end{conjecture}
The conjecture is certified by exact clique-LP certificates for $T=4$ ($3\le d\le10$, $k\le12$; $11\le d\le14$, $k\le10$), $T=5$
($3\le d\le9$, $k\le6$), all $d\le16$ with $T+k\le5$, and $d=3$ with $T\le60$ ($k=1$) and $T\le40$ ($k\le3$). On $d=3$ the margin
$1-\max\rho$ decays like $\approx2.9/T^2$, so a proof for all $T$ must be sharp in $T$.

\begin{table}[t]\centering
\caption{Torus values $\Ptor(4;T,k)$, computed from~\cite[Eq.~(16)]{Herasymenko2026} and independently from Lemma~\ref{lem:R}
(agreement $<10^{-9}$), and Werner's universal value $W(4,T,k)=D_T/D_{T+k}$.}\label{tab:torus}
\begin{tabular}{ccc}\toprule
$(T,k)$ & $\Ptor(4;T,k)$ & $W(4,T,k)$\\\midrule
(1,1)&0.437500&0.400000\\
(1,2)&0.281250&0.200000\\
(2,1)&0.577816&0.500000\\
(1,3)&0.234375&0.114286\\
(2,2)&0.428776&0.285714\\
(3,1)&0.682631&0.571429\\\bottomrule
\end{tabular}
\end{table}


\subsection{Negative results}\label{sec:neg}
Sidon alone does not suffice beyond $T+k=2$. For $K=\{18,21,27,28\}\subset\Z_{32}$ (Sidon, not $B_3$),
$P_K(1,1)=0.4375=\Ptor$ but $P_K(1,2)=0.291797>\Ptor=0.281250$ and $P_K(2,1)=0.607913>\Ptor=0.577816$.
\begin{table}[tbp]\centering\small
\caption{Counterfeiting values $P_K(T,k)$ at $T+k\in\{3,4\}$ for the sets of Table~\ref{tab:negsets}, from the reduced dual of Lemma~\ref{lem:R} (QH-Lab \texttt{qhlab.lemma\_r}). Bold: largest value in the row. A matched $B_{T+k}$ set of the same size gives $\Ptor$ in every row. All entries agree with the independent R5 sweep to $\le5\times10^{-8}$; for $d=4$ they also agree with a primal (lower-bound) SDP to $\le10^{-6}$.}\label{tab:neg}
\begin{tabular}{cc c cccc}\toprule
$d$ & $(T,k)$ & $\Ptor(d;T,k)$ & Sidon only & Singer & AP & random\\\midrule
4 & (1,2) & 0.281250 & 0.295129 & 0.343750 & \textbf{0.686077} & 0.434341\\
4 & (2,1) & 0.577816 & 0.624386 & 0.814044 & \textbf{0.918002} & 0.804521\\
4 & (1,3) & 0.234375 & 0.248486 & 0.331215 & \textbf{0.604402} & 0.409700\\
4 & (2,2) & 0.428776 & 0.519855 & 0.764730 & 0.875034 & \textbf{0.907654}\\
4 & (3,1) & 0.682631 & 0.800184 & \textbf{0.985848} & 0.969445 & 0.937383\\
\midrule
6 & (1,2) & 0.138889 & 0.168386 & 0.212963 & \textbf{0.674107} & 0.378659\\
6 & (2,1) & 0.421234 & 0.456406 & 0.715224 & \textbf{0.930640} & 0.776000\\
6 & (1,3) & 0.083333 & 0.141279 & 0.212560 & \textbf{0.600597} & 0.245909\\
6 & (2,2) & 0.232883 & 0.304882 & 0.674970 & \textbf{0.881735} & 0.669629\\
6 & (3,1) & 0.517533 & 0.714506 & \textbf{0.983498} & 0.976595 & 0.909976\\
\midrule
8 & (1,2) & 0.082031 & 0.119188 & 0.152344 & \textbf{0.669905} & 0.368231\\
8 & (2,1) & 0.330058 & 0.445114 & 0.666175 & \textbf{0.936120} & 0.851925\\
8 & (1,3) & 0.038086 & 0.095927 & 0.148619 & \textbf{0.599266} & 0.309036\\
8 & (2,2) & 0.144007 & 0.363933 & 0.628594 & \textbf{0.885440} & 0.800783\\
8 & (3,1) & 0.413747 & 0.792293 & 0.984480 & 0.979150 & \textbf{0.988168}\\
\bottomrule
\end{tabular}
\end{table}

\begin{table}[tbp]\centering\small
\caption{Key sets used in Table~\ref{tab:neg} (one set per $d$ and type, the same for every $(T,k)$ except Sidon-only and random, which were drawn per row in the R5 sweep). ``Sidon only'' sets are $B_2$ but not $B_3$; Singer sets are perfect difference sets, hence Sidon and not $B_3$; AP is $\{0,\dots,d-1\}$ with $q=2^{20}$ (no wrap-around).}\label{tab:negsets}
\begin{tabular}{cc ll}\toprule
$d$ & $(T,k)$ & type & $K\subset\Z_q$\\\midrule
4 & (1,2) & Sidon only & $\{13,20,21,31\}\subset\Z_{32}$\\
4 & all & Singer & $\{0,1,3,9\}\subset\Z_{13}$\\
4 & all & AP & $\{0,1,2,3\}\subset\Z_{1048576}$\\
4 & (1,2) & random & $\{1,7,8,11\}\subset\Z_{16}$\\
4 & (2,1) & Sidon only & $\{5,10,27,29\}\subset\Z_{32}$\\
4 & (2,1) & random & $\{6,8,9,11\}\subset\Z_{16}$\\
4 & (1,3) & Sidon only & $\{0,4,12,27\}\subset\Z_{32}$\\
4 & (1,3) & random & $\{6,8,9,11\}\subset\Z_{16}$\\
4 & (2,2) & Sidon only & $\{7,8,10,19\}\subset\Z_{32}$\\
4 & (2,2) & random & $\{0,10,12,14\}\subset\Z_{16}$\\
4 & (3,1) & Sidon only & $\{3,10,22,25\}\subset\Z_{32}$\\
4 & (3,1) & random (Sidon) & $\{0,7,10,12\}\subset\Z_{16}$\\
\midrule
6 & (1,2) & Sidon only & $\{1,29,50,110,116,117\}\subset\Z_{128}$\\
6 & all & Singer & $\{0,1,3,10,14,26\}\subset\Z_{31}$\\
6 & all & AP & $\{0,1,2,3,4,5\}\subset\Z_{1048576}$\\
6 & (1,2) & random & $\{0,11,15,17,21,29\}\subset\Z_{32}$\\
6 & (2,1) & Sidon only & $\{4,45,67,72,95,125\}\subset\Z_{128}$\\
6 & (2,1) & random & $\{1,5,6,17,29,30\}\subset\Z_{32}$\\
6 & (1,3) & Sidon only & $\{11,12,99,117,123,125\}\subset\Z_{128}$\\
6 & (1,3) & random & $\{5,9,11,27,28,31\}\subset\Z_{32}$\\
6 & (2,2) & Sidon only & $\{14,25,30,60,82,121\}\subset\Z_{128}$\\
6 & (2,2) & random & $\{2,4,5,23,25,29\}\subset\Z_{32}$\\
6 & (3,1) & Sidon only & $\{0,10,23,43,91,126\}\subset\Z_{128}$\\
6 & (3,1) & random & $\{5,9,13,15,22,31\}\subset\Z_{32}$\\
\midrule
8 & (1,2) & Sidon only & $\{51,53,59,83,93,104,121,122\}\subset\Z_{128}$\\
8 & all & Singer & $\{0,1,3,13,32,36,43,52\}\subset\Z_{57}$\\
8 & all & AP & $\{0,1,2,3,4,5,6,7\}\subset\Z_{1048576}$\\
8 & (1,2) & random & $\{0,6,8,12,16,21,24,27\}\subset\Z_{32}$\\
8 & (2,1) & Sidon only & $\{0,15,43,60,78,92,97,104\}\subset\Z_{128}$\\
8 & (2,1) & random & $\{1,2,9,10,14,16,17,29\}\subset\Z_{32}$\\
8 & (1,3) & Sidon only & $\{14,30,80,84,85,95,98,117\}\subset\Z_{128}$\\
8 & (1,3) & random & $\{1,5,7,10,16,17,23,30\}\subset\Z_{32}$\\
8 & (2,2) & Sidon only & $\{16,45,51,52,82,91,106,110\}\subset\Z_{128}$\\
8 & (2,2) & random & $\{0,2,6,16,21,27,28,31\}\subset\Z_{32}$\\
8 & (3,1) & Sidon only & $\{0,24,43,50,60,108,117,122\}\subset\Z_{128}$\\
8 & (3,1) & random & $\{1,6,8,13,17,21,27,31\}\subset\Z_{32}$\\
\bottomrule\end{tabular}
\end{table}

Table~\ref{tab:neg} extends this to $T+k\in\{3,4\}$ and $d\in\{4,6,8\}$. In all 15 rows every set that is Sidon but not $B_{T+k}$
(the Sidon-only sets and the Singer sets) has $P_K(T,k)>\Ptor$, and at fixed $T+k$ the margin grows with $T$. Singer sets, which are optimal
for bias and form equiangular tight frames, are always worse than the Sidon-only sets and come within $0.017$ of $1$ at $(3,1)$,
where they are the worst set tested for $d=4,6$. Arithmetic progressions, which have large bias, are the worst set in 11 of the
15 rows. Small Fourier bias therefore does not imply counterfeit resistance; what matters is the absence of
additive coincidences of order $T+k$.
These separations are numerical: the entries are reduced-dual values (for $d=4$ also bracketed by a primal SDP), not certificates.
The failure of Sidon-only sets is also not universal in $d$: for $d=2$ and $(T,k)=(1,2)$, every Sidon set that is not $B_3$ among the
24 we tested (e.g.\ $K=\{0,26\}\subset\Z_{39}$) attains $\Ptor(2;1,2)=0.75$, whereas for $d=2$, $(T,k)=(2,1)$ and for $d=3$ at
$T+k\in\{3,4\}$ none of the tested sets does (24 and 72 sets, $q<40$).

\subsection{The case $T=1$: an exact characterisation}\label{sec:t1iff}
Fix $k\ge1$. A \emph{charge} is a vector $c\in\Z^d$ with $\sum_ic_i=k$ and $c+e_a\ge0$ for some $a\in[d]$; let $\mathcal C_k$ be the
(finite) set of charges, and put
\[
M(c)=\sum_{a\in[d]}\mult(c+e_a)\qquad(\mult(\gamma):=0\text{ if }\gamma\not\ge0),\qquad M^\star_k=\max_{c\in\mathcal C_k}M(c).
\]
A charge with $M(c)=M^\star_k$ is \emph{tight}. Charges $c\ne c'$ are \emph{merged} by $K$ if $K\cdot c\equiv K\cdot c'\pmod q$, where
$K\cdot c=\sum_ic_ik_i$. Every charge has $M(c)\ge1$. For $s\in\Z_q$ let $\Lambda_K(s)=\sum_{a\in[d]}r_{k+1}(s+k_a)$.

\begin{lemma}[Aggregation]\label{lem:agg}
$\Lambda_K(s)=\sum_{c\in\mathcal C_k:\,K\cdot c\equiv s}M(c)$ for every $s\in\Z_q$.
\end{lemma}
\begin{proof}
$r_{k+1}(v)=\sum_{|\beta|=k+1,\,K\cdot\beta\equiv v}\mult(\beta)$. Substituting $c=\beta-e_a$,
$\sum_a\sum_{K\cdot\beta\equiv s+k_a}\mult(\beta)=\sum_a\sum_{c:\,c+e_a\ge0,\,K\cdot c\equiv s}\mult(c+e_a)=\sum_{c\in\mathcal C_k,\,K\cdot c\equiv s}M(c)$.
\end{proof}

\begin{lemma}[Tight charges]\label{lem:t1tight}
For all $d\ge2$, $k\ge1$, the tight charges in $\mathcal C_k$ are exactly the balanced vectors $c\ge0$ with $|c|=k$ and entries in
$\{\lfloor k/d\rfloor,\lfloor k/d\rfloor+1\}$; every other charge has $M(c)<M^\star_k$. In particular
$M^\star_k=\mult(c^\ast)\sum_a\frac{k+1}{c^\ast_a+1}$ for any balanced $c^\ast$.
\end{lemma}
\begin{proof}
\emph{Charges with a negative entry.} Such a charge is $c=\beta-e_a$ with $|\beta|=k+1$ and $\beta_a=0$, so $M(c)=\mult(\beta)$.
Pick $b$ with $\beta_b\ge1$. Then $c'=\beta-e_b\ge0$ is a charge, and
$M(c')\ge\mult(c'+e_b)+\mult(c'+e_a)=\mult(\beta)+\mult(\beta-e_b+e_a)>\mult(\beta)=M(c)$.

\emph{Nonnegative charges.} For $c\ge0$ with $|c|=k$ we have $M(c)=\mult(c)\sum_a\frac{k+1}{c_a+1}$. Suppose $c_i=x\ge y+2$ with
$y=c_j$, and put $c'=c-e_i+e_j$ and $R=\sum_{a\ne i,j}\frac1{c_a+1}\ge0$. Then $\mult(c')=\mult(c)\,\frac{x}{y+1}$, and
\[
\frac{M(c')}{(k+1)\mult(c)}=\frac1{y+1}+\frac{x}{(y+1)(y+2)}+\frac{x}{y+1}R\ >\ \frac1{y+1}+\frac1{x+1}+R=\frac{M(c)}{(k+1)\mult(c)} .
\]
This holds because $x/(y+1)>1$ and $x(x+1)>(y+1)(y+2)$. So any non-balanced $c\ge0$ is strictly beaten by a less unbalanced charge.
The balanced charges are permutations of one another, so they all have the same value.
\end{proof}

\begin{proposition}\label{prop:t1tor}
$\Ptor(d;1,k)=M^\star_k/d^{k+1}$. Moreover the program
\begin{equation}\label{eq:t1tor}
\min_{z\in(0,\infty)^d}\sum_a\frac1{z_a}\quad\text{s.t.}\quad d^{-(k+2)}\sum_a\mult(c+e_a)\,z_a\le1\quad\forall c\in\mathcal C_k
\end{equation}
has value $\Ptor(d;1,k)$ and the unique minimiser $z^\circ=(d^{k+2}/M^\star_k)\mathbf 1$.
\end{proposition}
\begin{proof}
The feasible set of \eqref{eq:t1tor} is convex and invariant under permutations of $[d]$ (they permute $\mathcal C_k$), and the objective
is strictly convex and symmetric. Hence the minimiser is unique (the midpoint of two minimisers would be strictly better) and
permutation invariant, $z=\zeta\mathbf 1$; such $z$ is feasible iff $\zeta M^\star_k\le d^{k+2}$, so the value is $M^\star_k/d^{k+1}$.
Now let $K'=\{1,b,b^2,\dots,b^{d-1}\}$ with $b=k+3$, and $q'>(k+2)b^{d-1}$. Sums of $(k+2)$-multisets from $K'$ have base-$b$ digits $\le k+2$,
so $K'$ is $B_{k+2}$ mod $q'$, and two charges $c=\beta-e_a$, $c'=\beta'-e_{a'}$ with $K'\cdot c\equiv K'\cdot c'$ satisfy
$K'\cdot(\beta+e_{a'})\equiv K'\cdot(\beta'+e_a)$, hence $\beta+e_{a'}=\beta'+e_a$ and $c=c'$. By Lemmas~\ref{lem:R} and~\ref{lem:agg}, $P_{K'}(1,k)$
is therefore exactly \eqref{eq:t1tor}. Since $K'$ is $B_{k+1}$, Theorem~\ref{thm:torus}(b) gives $P_{K'}(1,k)=\Ptor(d;1,k)$.
\end{proof}

\begin{theorem}[$T=1$ characterisation]\label{thm:t1iff}
For every $K\subset\Z_q$ with $|K|=d\ge2$ and every $k\ge1$,
\[
P_K(1,k)=\Ptor(d;1,k)\iff\max_{s\in\Z_q}\Lambda_K(s)\le M^\star_k ,
\]
and otherwise $P_K(1,k)>\Ptor(d;1,k)$. Equivalently, equality holds iff no class $\{c\in\mathcal C_k:K\cdot c\equiv s\}$ has total
weight $\sum M(c)$ above $M^\star_k$.
\end{theorem}
\begin{proof}
With $T=1$ and distinct keys, $r_1=\mathbf 1_K$, so by Lemmas~\ref{lem:R} and~\ref{lem:agg}, $P_K(1,k)$ is the minimum of $\sum_a1/z_a$ subject
to $d^{-(k+2)}\sum_{c:\,K\cdot c\equiv s}\sum_a\mult(c+e_a)z_a\le1$ for all $s$. Each such constraint is a sum of nonnegative terms, one per
charge in the class, so every $K$-feasible $z$ is feasible for \eqref{eq:t1tor}; this gives $P_K(1,k)\ge\Ptor$ (Lemma~\ref{lem:lower}). If
equality holds, the $K$-optimal $z$, which exists by Lemma~\ref{lem:R}, minimises \eqref{eq:t1tor}, so $z=z^\circ$ by
Proposition~\ref{prop:t1tor}. Then $K$-feasibility of $z^\circ$ reads $\Lambda_K(s)\le M^\star_k$ for all $s$. Conversely, if this holds,
$z^\circ$ is $K$-feasible with value $\Ptor$.
\end{proof}

\begin{proposition}[Lagrangian form of necessity]\label{prop:t1lag}
Let $\mathcal T$ be the set of tight charges and, for $\mu\ge0$, $\lambda_s=\mu\,\#\{t\in\mathcal T:K\cdot t\equiv s\}$. Then
\[
P_K(1,k)\ \ge\ 2\sum_a\sqrt{c_a}-\mu|\mathcal T|,\qquad c_a=d^{-(k+2)}\sum_s\lambda_sr_{k+1}(s+k_a)\ \ge\ c^\circ:=\frac{\mu|\mathcal T|M^\star_k}{d^{k+3}},
\]
with equality $c_a=c^\circ$ for all $a$ iff no tight charge is merged with another charge. At $\mu|\mathcal T|=M^\star_k/d^{k+1}$ the
right-hand side equals $\Ptor(d;1,k)$ when no tight charge is merged, and exceeds it otherwise.
\end{proposition}
\begin{proof}
The first inequality is the Lagrangian bound in the proof of Theorem~\ref{thm:sidon}, with $r_2$ replaced by $r_{k+1}$. By
Lemma~\ref{lem:agg} applied to the single index $a$, $r_{k+1}(K\cdot t+k_a)=\sum_{c:\,K\cdot c\equiv K\cdot t}\mult(c+e_a)\ge\mult(t+e_a)$, with
strict inequality iff the class of $t$ contains some $c\ne t$ with $c+e_a\ge0$. So
$c_a\ge d^{-(k+2)}\mu\sum_{t\in\mathcal T}\mult(t+e_a)$. The last sum does not depend on $a$, because $\mathcal T$ is permutation invariant,
and summing over $a$ gives $|\mathcal T|M^\star_k$, so it equals $|\mathcal T|M^\star_k/d$. If a tight $t$ is merged with some $c\ne t$,
pick $a$ with $c+e_a\ge0$; then $c_a>c^\circ$. Finally, with $u=\mu|\mathcal T|$, $2d\sqrt{c^\circ}-u$ is maximised at $u=M^\star_k/d^{k+1}$,
where it equals $M^\star_k/d^{k+1}=\Ptor$.
\end{proof}

\begin{corollary}[Necessity of the merging condition]\label{cor:t1nec}
If some tight charge is merged with another charge, then $P_K(1,k)>\Ptor(d;1,k)$.
\end{corollary}
\begin{proof}
The class of that tight charge has weight at least $M^\star_k+1$; apply Theorem~\ref{thm:t1iff}, or Proposition~\ref{prop:t1lag} for an
explicit lower bound.
\end{proof}

Say that $K$ satisfies condition $(\mathrm H_k)$ if no tight charge is merged with another charge. By Corollary~\ref{cor:t1nec},
$(\mathrm H_k)$ is necessary for $P_K(1,k)=\Ptor$. It is also sufficient in the following cases.

\begin{lemma}\label{lem:binom}
Let $N\ge2$, $m=\lfloor N/2\rfloor$, and $u,v\ge0$ with $u+v\le m-1$. Then $\binom Nu+\binom Nv\le\binom Nm$.
\end{lemma}
\begin{proof}
Assume $u\le v$ and put $t=m-1-u$, so $u\le v\le t$ and $\binom Nv\le\binom Nt$. Further,
$\binom N{t+1}-\binom Nt=\binom Nt\frac{N-2t-1}{t+1}\ge\binom Nt\frac{2u+1}{m-u}$, since $N-2t-1\ge2u+1$ and $t+1=m-u$. If $t\ge u+1$, then
$\binom Nt\ge\binom N{u+1}=\binom Nu\frac{N-u}{u+1}$, and $\frac{(2u+1)(N-u)}{(u+1)(m-u)}\ge1$. If $t=u$, then $m=2u+1$ and
$\frac{2u+1}{m-u}\ge1$. In both cases $\binom N{t+1}-\binom Nt\ge\binom Nu$, so
$\binom Nu+\binom Nv\le\binom N{t+1}\le\binom Nm$, because $t+1\le m$.
\end{proof}

\begin{theorem}[Sufficiency of $(\mathrm H_k)$ in three regimes]\label{thm:t1suff}
$P_K(1,k)=\Ptor(d;1,k)$ iff $(\mathrm H_k)$ holds, in each of the following cases.
\begin{enumerate}
\item[(a)] $k=1$, all $d$: here $(\mathrm H_1)$ is the Sidon property (Theorem~\ref{thm:sidon}).
\item[(b)] $d=2$, all $k$: here $(\mathrm H_k)$ holds iff $n:=\mathrm{ord}(k_2-k_1)\ge\lceil k/2\rceil+2$.
\item[(c)] $k=2$, all $d$: here $(\mathrm H_2)$ holds iff the only coincidences among sums of $3$-multisets from $K$ are
$3k_i\equiv3k_l$, i.e.\ every $3$-multiset with two distinct elements has a unique sum.
\end{enumerate}
\end{theorem}
\begin{proof}
Necessity is Corollary~\ref{cor:t1nec}; we prove sufficiency through Theorem~\ref{thm:t1iff}.

(a) For $k=1$, $M(e_j)=2d-1$, and every other charge $c=e_b+e_{b'}-e_a$ has $M(c)\le2<2d-1$. If $k_i+k_j\equiv k_l+k_m$ with
$\{i,j\}\ne\{l,m\}$, then the tight charge $e_i$ is merged with $e_l+e_m-e_j\ne e_i$, so $(\mathrm H_1)$ implies Sidon. For Sidon $K$, the
dual part of the proof of Theorem~\ref{thm:sidon} shows that $\Lambda_K\le2d-1$.

(b) Let $N=k+2$ and $g=k_2-k_1$. The charges are $c=(k-j,j)$, $j=-1,\dots,k+1$, with
$M(c)=\binom{k+1}j+\binom{k+1}{j+1}=\binom N{i}$, where $i=j+1\in[0,N]$, and $K\cdot c=kk_1+jg$. So two charges are merged iff their indices
$i$ agree mod $n$. The tight indices are $\lfloor N/2\rfloor$ and $\lceil N/2\rceil$, so $(\mathrm H_k)$ holds iff $n>\lceil N/2\rceil$,
i.e.\ $n\ge\lceil k/2\rceil+2$. Then every class has at most two indices $i<i+n$, and with $v=N-i-n$ we get
$i+v=N-n\le\lfloor N/2\rfloor-1$. Lemma~\ref{lem:binom} gives $\binom Ni+\binom N{i+n}=\binom Ni+\binom Nv\le\binom N{\lfloor N/2\rfloor}=M^\star_k$.

(c) Let $d\ge3$ ($d=2$ is covered by (b)). The weights are $M(e_i+e_j)=6d-6$ for $i\ne j$, $M(2e_i)=3d-2$, and
$M(\beta-e_a)=\mult(\beta)\in\{1,3,6\}$ for $|\beta|=3$, $\beta_a=0$. So the tight charges are exactly the $e_i+e_j$, $i\ne j$.
A tight $e_i+e_j$ is merged with $c'=\beta-e_a$ ($\beta\ge0$) iff $K\cdot(e_i+e_j+e_a)\equiv K\cdot\beta$ with $\beta\ne e_i+e_j+e_a$. Every
$3$-multiset with two distinct elements has the form $\{i,j,a\}$ with $i\ne j$, so this proves the stated form of $(\mathrm H_2)$.

Assume $(\mathrm H_2)$ and fix a class $S$; we show $\sum_{c\in S}M(c)\le6d-6$.
\begin{itemize}
\item If $S$ contains a tight charge, $S$ is that charge alone.
\item If $S\ni2e_i$, consider any other $c'\in S$. It is not tight. It is not $2e_l$: from $2k_i\equiv2k_l$ we would get
$\{i,i,x\}\sim\{l,l,x\}$, contradicting $(\mathrm H_2)$. So $c'=\beta-e_a$ with $\beta_a=0$. Then $2e_i+e_a$ and $\beta$ are distinct
$3$-multisets with equal sums, so both are pure, which forces $a=i$ and $\beta=3e_l$. Hence $\sum_SM\le3d-2+(d-1)=4d-3\le6d-6$.
\item Otherwise every member has the form $c_\nu=\beta_\nu-e_{a_\nu}$ with $\beta_\nu(a_\nu)=0$, and this representation is unique. We use
three facts.
\begin{enumerate}
\item[(1)] If $a_\nu=a_\mu$ for $\nu\ne\mu$, then $K\cdot\beta_\nu\equiv K\cdot\beta_\mu$, so both $\beta$'s are pure.
\item[(2)] If $\beta_\nu$ has three distinct elements, then for every $\mu\ne\nu$ the $4$-multisets $\beta_\nu+e_{a_\mu}$ and
$\beta_\mu+e_{a_\nu}$ are disjoint. They have equal sums and are distinct. A common element $x$ could be removed, and this would
leave distinct $3$-multisets with equal sums; the first of these has at least two distinct elements, contradicting $(\mathrm H_2)$. So the
supports of $\beta_\nu$ and $\beta_\mu$ are disjoint.
\item[(3)] For each $l$ there is at most one member $3e_l-e_a$, since $a$ is determined by $k_a\equiv3k_l-s$.
\end{enumerate}
Let $x,y,w$ be the numbers of members with $\mult(\beta)=6,3,1$, and let $U$ be the union of the supports of the type-6 $\beta$'s; by (2)
these supports are pairwise disjoint, so $|U|=3x\le d$. By (1), two members with the same index $a$ both have pure $\beta$'s, so the type-6
and type-3 members have pairwise distinct indices $a_\nu$, and $x+y\le d$. For a type-1 member $\mu$, $\beta_\mu=3e_l$ has support $\{l\}$, so (2), applied to each type-6 member $\nu$ and this $\mu$, gives
$l\notin\operatorname{supp}\beta_\nu$; hence $l\notin U$. By (3) these $l$ are pairwise distinct, so $w\le d-3x$. Hence $\sum_SM=6x+3y+w\le6x+3(d-x)+(d-3x)=4d\le6d-6$, since $d\ge3$.
\end{itemize}
\end{proof}

\begin{conjecture}\label{conj:t1H}
For all $d\ge2$ and $k\ge1$: $P_K(1,k)=\Ptor(d;1,k)$ iff $(\mathrm H_k)$ holds.
\end{conjecture}

\begin{remark}[The $d=2$ exception and $d\ge3$]
Theorem~\ref{thm:t1suff} explains the $d=2$, $(T,k)=(1,2)$ observation of Section~\ref{sec:neg}. For $d=2$, $B_3$ means $n\ge4$, but
$P_K(1,2)=\Ptor$ already holds iff $n\ge3$. Thus $\{0,26\}\subset\Z_{39}$ ($n=3$) attains $0.75$, and in general $d=2$ sets with
$\lceil k/2\rceil+2\le n\le k+1$ attain $\Ptor(2;1,k)$ without being $B_{k+1}$. The mechanism is not specific to $d=2$: by (c), a Sidon
set whose only $3$-relations are $3k_i\equiv3k_l$ attains $\Ptor(d;1,2)$. Examples are $\{0,1,13\}\subset\Z_{39}$ ($d=3$, value $4/9$),
$\{0,10,11,14\}\subset\Z_{30}$ ($d=4$, $9/32$) and $\{0,23,34,36,50\}\subset\Z_{69}$ ($d=5$, $24/125$), none of them $B_3$. Similarly
$\{0,10,26,29\}\subset\Z_{30}$ attains $\Ptor(4;1,3)=15/64$ without being $B_3$. So $B_{T+k}$ is \emph{not} necessary for
$P_K=\Ptor$ when $d\ge3$, at least for $T=1$. The Sidon-only sets in the tables of Section~\ref{sec:neg} fail because they carry a non-pure $3$-relation; for
example $\{18,21,27,28\}\subset\Z_{32}$ has $2\cdot18+27\equiv3\cdot21$, a tight charge is merged, and $\max_s\Lambda_K(s)=19>18$.
\end{remark}

\begin{remark}[No finite codebook works for all $k$]
Since $\sum_s\Lambda_K(s)=d^{k+2}$, we have $\max_s\Lambda_K(s)\ge d^{k+2}/q$. On the other hand,
$M^\star_k/d^{k+2}\le\max_\beta\mult(\beta)/d^{k+1}\to0$ as $k\to\infty$, since this is the largest point probability of a uniform
multinomial. Hence for every finite $K\subset\Z_q$, $P_K(1,k)>\Ptor(d;1,k)$ for all sufficiently large $k$; precisely, this holds as soon as
$\Ptor(d;1,k)<d/q$.
\end{remark}

\section{Fusion and additive energy}\label{sec:fusion}
\begin{theorem}[Fusion]\label{thm:fusion}
For $x,y$ independent and uniform on $\Z_q$, the optimal average fidelity of a channel $\psi_x\otimes\psi_y\mapsto\psi_{x+y}$ is
$F_{\rm fuse}(K)=E_q(K)/d^3$. In particular $F_{\rm fuse}(K)\ge(2d-1)/d^2$, with equality iff $K$ is Sidon.
\end{theorem}
\begin{proof}
Let $r(v)=\#\{(a,b):k_b-k_a\equiv v\}$, so $\sum_v r(v)^2=E_q(K)$, and for each $v$ let $S_v=\{j:k_j+v\in K\}$, $|S_v|=r(v)$, with
$l(j,v)$ the index of $k_j+v$.

\emph{Achievability.} Define $A_v=\sum_{j\in S_v}\ket j\bra{j,l(j,v)}$. Each basis pair $\ket{j,l}$ lies in exactly one branch
($v=k_l-k_j$), so $\sum_vA_v^\dagger A_v=I$ and $\{A_v\}$ is a channel. On the input,
\[
A_v(\psi_x\otimes\psi_y)=\frac1d\sum_{j\in S_v}\omega^{xk_j+y(k_j+v)}\ket j=\frac{\omega^{yv}}d\sum_{j\in S_v}\omega^{(x+y)k_j}\ket j .
\]
The $y$-dependent factor is a global phase of the branch, so $\big|\bra{\psi_{x+y}}A_v(\psi_x\otimes\psi_y)\big|^2=r(v)^2/d^3$, and the
output fidelity is $\sum_v r(v)^2/d^3=E_q(K)/d^3$ for \emph{every} $(x,y)$.

\emph{Optimality.} The Choi operator of the target ensemble is
$X=\E_{x,y}\,\bar\psi_x\otimes\bar\psi_y\otimes\psi_{x+y}$, with entries
$X_{(jlm),(j'l'm')}=d^{-3}\,[k_m-k_j\equiv k_{m'}-k_{j'}]\,[k_m-k_l\equiv k_{m'}-k_{l'}]$. Hence $X=d^{-3}\sum_{(s,t)}\ket{\chi_{st}}\bra{\chi_{st}}$,
where $\chi_{st}$ is the 0/1 indicator of $B_{st}=\{(j,l,m):k_m-k_j\equiv s,\ k_m-k_l\equiv t\}$; the blocks $B_{st}$ partition the basis,
and on $B_{st}$ the difference $k_l-k_j\equiv s-t=:v$ is constant. Let $Y=\mathrm{diag}_{(j,l)}\,r(k_l-k_j)/d^3$. Since $Y\otimes I$ is
diagonal and positive on each block, $Y\otimes I\succeq d^{-3}\ket{\chi_{st}}\bra{\chi_{st}}$ iff $\sum_{B_{st}}d^{-3}/(r(v)/d^3)\le1$, i.e.\
$|B_{st}|\le r(v)$. The map $(j,l,m)\mapsto(l,j)$ sends $B_{st}$ injectively into $\{(a,b):k_b-k_a\equiv -v\}$, whose size is
$r(-v)=r(v)$. So $Y$ is dual feasible, and $\tr Y=\sum_{(j,l)}r(k_l-k_j)/d^3=\sum_vr(v)^2/d^3=E_q(K)/d^3$.

Finally $E_q(K)\ge2d^2-d$ (the trivial solutions $a=b,c=e$ and $a=c,b=e$), with equality iff $K$ is Sidon, giving the corollary.
\end{proof}

For $m\ge2$ let $r(v)=\#\{(a,b)\in[d]^2:k_b-k_a\equiv v\}$ as in the proof of Theorem~\ref{thm:fusion}, and define the
\emph{higher energy}
\[
E_m(K)=\sum_{v\in\Z_q}r(v)^m ,
\]
so that $E_2(K)=E_q(K)$. For $j=(j_1,\dots,j_m)\in[d]^m$ put $\delta(j)=(k_{j_2}-k_{j_1},\dots,k_{j_m}-k_{j_1})\in\Z_q^{m-1}$ and,
for $\delta=(\delta_2,\dots,\delta_m)\in\Z_q^{m-1}$ (with $\delta_1:=0$),
\[
N(\delta)=\#\{j\in[d]^m:\delta(j)=\delta\}=\#\{a\in[d]:k_a+\delta_i\in K\ \text{for all } i\},
\]
the second form because the keys are distinct, so $j$ is determined by $j_1$ and $\delta(j)$.

\begin{lemma}\label{lem:Em}
$\sum_{\delta\in\Z_q^{m-1}}N(\delta)^2=E_m(K)$.
\end{lemma}
\begin{proof}
$\sum_\delta N(\delta)^2$ counts triples $(a,b,\delta)$ with $k_a+\delta_i\in K$ and $k_b+\delta_i\in K$ for $i=2,\dots,m$. For fixed $(a,b)$ the
coordinates $\delta_i$ are chosen independently from $(K-k_a)\cap(K-k_b)$, a set of size $\#\{c\in K:c-k_a+k_b\in K\}=r(k_b-k_a)$.
Hence the count is $\sum_{a,b}r(k_b-k_a)^{m-1}=\sum_v r(v)\cdot r(v)^{m-1}=E_m(K)$.
\end{proof}

\begin{theorem}[$m$-fold fusion]\label{thm:fusionm}
Let $m\ge2$ and let $x_1,\dots,x_m$ be independent and uniform on $\Z_q$. The optimal average fidelity of a channel
$(\C^d)^{\otimes m}\to\C^d$ for $\psi_{x_1}\otimes\dots\otimes\psi_{x_m}\mapsto\psi_{x_1+\dots+x_m}$ is
\[
F_m(K)=\frac{E_m(K)}{d^{m+1}} ,
\]
and it is attained by a channel whose fidelity equals $F_m(K)$ on \emph{every} input $(x_1,\dots,x_m)$. For $m=2$ this is
Theorem~\ref{thm:fusion}.
\end{theorem}
\begin{proof}
Write $X_\Sigma=x_1+\dots+x_m$.

\emph{Achievability.} For $\delta\in\Z_q^{m-1}$ let $A_\delta=\sum_{j:\,\delta(j)=\delta}\ket{j_1}\bra{j}$. Every basis vector $\ket j$ of
$(\C^d)^{\otimes m}$ lies in exactly one branch, and distinct $j$ in one branch have distinct $j_1$, so
$\sum_\delta A_\delta^\dagger A_\delta=I$ and $\{A_\delta\}$ is a channel. Since $k_{j_i}=k_{j_1}+\delta_i$ on the branch,
\[
A_\delta\big(\psi_{x_1}\otimes\dots\otimes\psi_{x_m}\big)=d^{-m/2}\sum_{j:\,\delta(j)=\delta}\omega^{\sum_ix_ik_{j_i}}\ket{j_1}
=d^{-m/2}\,\omega^{\sum_ix_i\delta_i}\sum_{j:\,\delta(j)=\delta}\omega^{X_\Sigma k_{j_1}}\ket{j_1} .
\]
The factor $\omega^{\sum_ix_i\delta_i}$ is a global phase of the branch, so
$\big|\bra{\psi_{X_\Sigma}}A_\delta(\psi_{x_1}\otimes\dots\otimes\psi_{x_m})\big|^2=N(\delta)^2/d^{m+1}$, and the fidelity is
$\sum_\delta N(\delta)^2/d^{m+1}=E_m(K)/d^{m+1}$ for every input (Lemma~\ref{lem:Em}).

\emph{Optimality.} The average fidelity of a channel with Choi operator $J$ ($\tr_{\rm out}J=I$) is $\tr(JX)$ with
$X=\E_{x}\,\bar\psi_{x_1}\otimes\dots\otimes\bar\psi_{x_m}\otimes\psi_{X_\Sigma}$, and $\tr(JX)\le\tr\big(J(Y\otimes I)\big)=\tr Y$
whenever $Y\otimes I\succeq X$. Averaging $\omega^{x_i t}$ over $x_i$ gives
\[
X_{(j,l),(j',l')}=d^{-(m+1)}\prod_{i=1}^m\big[k_l-k_{j_i}\equiv k_{l'}-k_{j'_i}\big] .
\]
With $\sigma(j,l)=(k_l-k_{j_1},\dots,k_l-k_{j_m})\in\Z_q^m$ and $B_\sigma=\{(j,l):\sigma(j,l)=\sigma\}$, the blocks $B_\sigma$
partition the basis and $X=d^{-(m+1)}\sum_\sigma\ket{\chi_\sigma}\bra{\chi_\sigma}$, where $\chi_\sigma$ is the $0/1$ indicator of
$B_\sigma$. On $B_\sigma$ the input difference vector is constant, $\delta(j)=(\sigma_1-\sigma_i)_{i=2}^m=:\delta_\sigma$, and
\[
|B_\sigma|=\#\{l:k_l-\sigma_i\in K\ \forall i\}=\Big|\bigcap_i(K+\sigma_i)\Big|=\Big|\bigcap_i(K-(\delta_\sigma)_i)\Big|=N(\delta_\sigma),
\]
after translating by $-\sigma_1$. Let $Y=\mathrm{diag}_j\,N(\delta(j))/d^{m+1}$, which is positive since $N(\delta(j))\ge1$. Because $Y\otimes I$
is diagonal and $X$ is a direct sum over the blocks, $Y\otimes I\succeq X$ holds iff on every block
$\sum_{(j,l)\in B_\sigma}d^{-(m+1)}\big/\big(N(\delta_\sigma)/d^{m+1}\big)\le1$, i.e.\ $|B_\sigma|\le N(\delta_\sigma)$, which holds with
equality. Hence $Y$ is dual feasible and $\tr Y=\sum_jN(\delta(j))/d^{m+1}=\sum_\delta N(\delta)^2/d^{m+1}=E_m(K)/d^{m+1}$.
\end{proof}

\begin{corollary}\label{cor:fusionm}
For $m\ge2$, $E_m(K)\ge d^m+d^2-d$, with equality if and only if $K$ is a Sidon set. Hence
$F_m(K)\ge(d^m+d^2-d)/d^{m+1}$, with equality iff $K$ is Sidon. Moreover $F_m(K)\le1$, with equality iff $K$ is a coset of a
subgroup of $\Z_q$ of order $d$.
\end{corollary}
\begin{proof}
$r(0)=d$ and $\sum_{v\ne0}r(v)=d^2-d$. For integers $r\ge0$ and $m\ge2$, $r^m\ge r$ with equality iff $r\in\{0,1\}$, so
$E_m(K)=d^m+\sum_{v\ne0}r(v)^m\ge d^m+d^2-d$, with equality iff $r(v)\le1$ for all $v\ne0$. This is the Sidon property: if
$k_a+k_e\equiv k_b+k_c$ with $\{a,e\}\ne\{b,c\}$, then the ordered pairs $(b,a)$ and $(e,c)$ are distinct and have the same nonzero
difference $k_a-k_b\equiv k_c-k_e$; conversely two distinct ordered pairs $(a,b)\ne(c,e)$ with $k_b-k_a\equiv k_e-k_c\not\equiv0$ give
$k_b+k_c\equiv k_a+k_e$ with $\{b,c\}\ne\{a,e\}$. For the upper bound, $r(v)\le d$ gives $E_m(K)\le d^{m-1}\sum_vr(v)=d^{m+1}$, with
equality iff $r(v)\in\{0,d\}$ for all $v$, i.e.\ $K+v=K$ for every $v\in K-K$, i.e.\ $K$ is a coset of a subgroup of order $d$.
\end{proof}

\begin{remark}
For $m\ge3$ the optimal measurement reads out the full difference vector $\delta\in\Z_q^{m-1}$, and the optimum is the
\emph{common-base-point} count $\sum_\delta N(\delta)^2$; Lemma~\ref{lem:Em} identifies it with $\sum_v r(v)^m$. For $m=3$ and
$K=\{0,1,3\}\subset\Z_7$ (Sidon) this gives $33/81$, for $\{0,1,2\}\subset\Z_7$ it gives $45/81$, and for $\{0,1,2,4\}\subset\Z_{11}$
it gives $100/256$. In the notation of~\cite{SchoenShkredov2013,Shkredov2013HigherEnergies}, $r(v)=(K\circ K)(v)=|K\cap(K-v)|$, $E_m(K)$ is the
higher energy $\mathsf E_m(K)$, and Lemma~\ref{lem:Em} is the identity $\mathsf E_m(A)=\sum_{s\in G^{m-1}}|A_s|^2$,
$A_s=A\cap(A-s_1)\cap\dots\cap(A-s_{m-1})$; see also the survey~\cite{Shkredov2015Survey}.
\end{remark}


Thus the same additive structure that makes a codebook hard to clone makes it hard to fuse: Sidon sets minimise both, and
perfect fusion forces orthogonal (hence clonable) states. Phase-state tokens are also malleable (the public shift $Z^y$ maps
$\psi_x\mapsto\psi_{x+y}$), and verification oracles leak the key~\cite{Nagaj2016}.

\section{Finite orbits of non-abelian groups}\label{sec:nonab}
Let $H$ be a compact group acting unitarily on $V$, $G\le H$ closed, and $X=G\cdot f\subset H\cdot f$ with uniform weights. Because
$X\subset H\cdot f$ and Haar measure is invariant, the $H$-twirl of $R_X$ is $R_H$, i.e.\ $\int_H h\,D\,h^{-1}dh=0$; this identity is used
in parts (b) and (e) below. Let
$u_x=\bar x^{(T)}\otimes x^{(T+k)}$, $R_X=\E_X u_xu_x^\dagger$, $R_H=\int_H R_{hf}\,dh$, $D=R_X-R_H$, and $P_\bullet=\min\{\tr Y:Y\otimes I\succeq R_\bullet\}$.

A natural guess, suggested by Proposition~\ref{prop:nearW} and by the torus case, is that matching the $(T+k)$-th moment,
$\E_X x^{\otimes(T+k)}=\E_H x^{\otimes(T+k)}$, implies $P_X=P_H$. It holds for unitary designs (Werner) and for orbits of
highest-weight (coherent) vectors~\cite{ChiribellaYang2014,CwiklinskiStudzinski2026}, where Lemma~\ref{lem:B} is tight. It fails in general:

\begin{theorem}[Counterexamples]\label{thm:counter}
(i) $H=SU(3)$ acting on $V=\Sym^2\C^3$ ($\dim V=6$), $G$ the Hessian group (order 216 mod phases), $(T,k)=(1,1)$: there is a
fiducial $f$ with $M_2(X)=M_2(H)$ exactly and $P_X-P_H\ge8.0399\times10^{-4}$.
(ii) $H=SU(2)$ acting on spin 1 (through $SO(3)$, i.e.\ the vector representation on $\C^3$), $G$ the octahedral rotation group (order 24), $(T,k)=(1,1)$: there is a fiducial with
$M_2(X)=M_2(H)$ exactly and $P_X-P_H\ge6.437\times10^{-4}$.
\end{theorem}
Both are certified in ball arithmetic, with exact moment matching established by a Poincar\'e--Miranda (resp.\ intermediate value)
argument; see Appendix~\ref{app:ca}. Independently re-implemented checks (exact Schur--Weyl twirl for $SU(3)$, exact Euler product
rule for $SO(3)$) with other fiducials give strictly positive gaps in every tuned case (4/4 and 4/4).

\begin{theorem}[Window theorem]\label{thm:window}
Assume (C1) $G$ and $H$ have the same commutant on $\Sym^TV$. Let $\mathcal Y_H$ be the $H$-invariant optimal duals for $R_H$; for
$Y\in\mathcal Y_H$ let $Z_Y=Y\otimes I-R_H\succeq0$ and $\Pi_Y$ the projector onto $\ker Z_Y$.
(a) $P_X\ge P_H$, with equality iff some $Y\in\mathcal Y_H$ has $Z_Y-D\succeq0$.
(b) If every $G$-invariant operator on $\ker Z_Y$ is $H$-invariant, then $\Pi_YD\Pi_Y=0$.
(c) If the hypothesis of (b) holds for every $Y\in\mathcal Y_H$ and $(1-\Pi_Y)D\Pi_Y\ne0$ for every $Y\in\mathcal Y_H$, then
$P_X>P_H$. (Proof idea: $\Pi_Y(Z_Y-D)\Pi_Y=0$, and a positive semidefinite operator with a vanishing diagonal block has vanishing
off-diagonal blocks; so no optimal $H$-dual is feasible for $X$, and (a) applies.)
(d) If some $Y$ has $\Pi_YD\Pi_Y=0$, $(1-\Pi_Y)D\Pi_Y=0$ and $\lVert(1-\Pi_Y)D(1-\Pi_Y)\rVert\le\lambda^+_{\min}(Z_Y)$, then $P_X=P_H$.
(e) (Selection rule) $(1-\Pi_Y)D\Pi_Y$ can only receive contributions from $H$-isotypic components $\lambda\ne\mathbf1$ of the moment
difference $M_{2T+k}(X)-M_{2T+k}(H)$ with $V_\lambda^G\ne0$ and $V_\lambda\hookrightarrow\mathrm{Hom}(\ker Z_Y,(\ker Z_Y)^\perp)$.
\end{theorem}
\paragraph{Actions.} Fix an orthonormal basis of $V$, write $\bar A$ for the entrywise complex conjugate and $A^{\rm t}$ for the transpose in
this basis, and put $n=2T+k$. Each $h\in H$ acts on
\begin{itemize}
\item[(A1)] $V$ by $h$, and on $\bar V$ (the same vector space with the conjugate representation) by $\bar h$. Since $h$ is unitary,
$\bar h=(h^{-1})^{\rm t}$, so $\bar V\cong V^\ast$ as $H$-modules;
\item[(A2)] $V^{\otimes n}$ by $h^{\otimes n}$; on $\bar V^{\otimes T}\otimes V^{\otimes(T+k)}$ by $\bar h^{\otimes T}\otimes h^{\otimes(T+k)}$; on the invariant
subspaces $\Sym^T(\bar V)$, $\Sym^{T+k}(V)$ and $\mathcal K=\Sym^T(\bar V)\otimes\Sym^{T+k}(V)$ by restriction. On $\mathcal K$ we write
$U_h=\bar h^{\otimes T}\otimes h^{\otimes(T+k)}|_{\mathcal K}$, and on the input space $\Sym^T(\bar V)$ of $Y$ we write $\bar h_T=\bar h^{\otimes T}|_{\Sym^T(\bar V)}$;
\item[(A3)] operators by conjugation: $\mathrm{Ad}_h(A)=h^{\otimes n}Ah^{\dagger\otimes n}$ on $\mathrm{End}(V^{\otimes n})$, $A\mapsto U_hAU_h^\dagger$ on
$\mathrm{End}(\mathcal K)$, $Y\mapsto\bar h_TY\bar h_T^\dagger$ on $\mathrm{End}(\Sym^T\bar V)$, and, for $H$-invariant subspaces $\mathcal A,\mathcal B\subset\mathcal K$,
$S\mapsto U_h|_{\mathcal B}\,S\,U_h|_{\mathcal A}^\dagger$ on $\mathrm{Hom}(\mathcal A,\mathcal B)$; as $H$-modules $\mathrm{Hom}(\mathcal A,\mathcal B)\cong\mathcal B\otimes\mathcal A^\ast$.
\end{itemize}
The same formulas with $g\in G$ define the $G$-actions. An operator is \emph{$G$-invariant} (resp.\ \emph{$H$-invariant}) if it is fixed by all
$g\in G$ (resp.\ $h\in H$). With $u_x=\bar x^{\otimes T}\otimes x^{\otimes(T+k)}\in\mathcal K$ we have $U_hu_x=u_{hx}$ and
$U_h(Y\otimes I)U_h^\dagger=(\bar h_TY\bar h_T^\dagger)\otimes I$. $R_X$ is $G$-invariant because $X$ is a $G$-orbit with uniform weights, and $R_H$ is
$H$-invariant by invariance of Haar measure. For every $x=gf\in X$, right invariance of Haar measure gives
$\int_HU_hu_xu_x^\dagger U_h^\dagger\,dh=\int_Hu_{hgf}u_{hgf}^\dagger\,dh=R_H$; averaging over $X$,
\begin{equation}\label{eq:W-twirl}
\int_HU_hR_XU_h^\dagger\,dh=R_H,\qquad\text{hence}\qquad\int_HU_hDU_h^\dagger\,dh=0 .
\end{equation}
The minima defining $P_X$ and $P_H$ are attained: a feasible $Y$ satisfies $Y\succeq0$, and $\{Y\succeq0:\tr Y\le c\}$ is compact.
For $Y\in\mathcal Y_H$, $Z_Y=Y\otimes I-R_H$ is $H$-invariant, hence $\ker Z_Y$ and $(\ker Z_Y)^\perp$ are $H$-invariant subspaces of $\mathcal K$,
$\Pi_Y$ commutes with every $U_h$, and $Z_Y=(1-\Pi_Y)Z_Y(1-\Pi_Y)$. Finally $Z_Y-D=Y\otimes I-R_X$, and $D$ is Hermitian.

\begin{lemma}[Zero diagonal block]\label{lem:zeroblock}
Let $M\succeq0$ on a finite-dimensional Hilbert space and let $\Pi$ be an orthogonal projection with $\Pi M\Pi=0$. Then $M\Pi=0$ and $\Pi M=0$; in
particular $(1-\Pi)M\Pi=0$.
\end{lemma}
\begin{proof}
For every vector $v$,
\[0=\langle\Pi v,M\Pi v\rangle=\lVert M^{1/2}\Pi v\rVert^2 .\]
So $M^{1/2}\Pi=0$, hence $M\Pi=M^{1/2}(M^{1/2}\Pi)=0$ and $\Pi M=(M\Pi)^\dagger=0$.
\end{proof}

\begin{proof}[Proof of Theorem~\ref{thm:window}]
\emph{(a)} Let $Y$ attain $P_X$ and put $Y_G=\int_G\bar g_TY\bar g_T^\dagger\,dg$ (normalised Haar measure on $G$; the uniform average if $G$ is
finite). Conjugating $Y\otimes I\succeq R_X$ by $U_g$ and using $G$-invariance of $R_X$ gives $(\bar g_TY\bar g_T^\dagger)\otimes I\succeq R_X$ for every $g$, so
$Y_G$ is feasible for $X$, and $\tr Y_G=\tr Y=P_X$. $Y_G$ commutes with every $\bar g_T$; taking entrywise conjugates, $\overline{Y_G}$ commutes with every
$g^{\otimes T}|_{\Sym^TV}$, so by (C1) it commutes with every $h^{\otimes T}|_{\Sym^TV}$, i.e.\ $Y_G$ commutes with every $\bar h_T$: $Y_G$ is
$H$-invariant. Conjugating $Y_G\otimes I\succeq R_X$ by $U_h$ gives $Y_G\otimes I\succeq U_hR_XU_h^\dagger$ for every $h$; integrating and using
\eqref{eq:W-twirl}, $Y_G\otimes I\succeq R_H$. Hence $P_H\le\tr Y_G=P_X$. If $P_X=P_H$, then $Y_G$ is an $H$-invariant optimal dual for $R_H$, i.e.\
$Y_G\in\mathcal Y_H$, and $Z_{Y_G}-D=Y_G\otimes I-R_X\succeq0$. Conversely, if $Y\in\mathcal Y_H$ and $Z_Y-D=Y\otimes I-R_X\succeq0$, then $Y$ is
feasible for $X$ and $P_X\le\tr Y=P_H$; with $P_X\ge P_H$ this gives equality.

\emph{(b)} The compression $B:=\Pi_YD\Pi_Y|_{\ker Z_Y}\in\mathrm{End}(\ker Z_Y)$ is $G$-invariant, because $D=R_X-R_H$ is $G$-invariant and $\Pi_Y$
commutes with every $U_g$. By hypothesis $B$ is $H$-invariant, so it equals its $H$-twirl:
$\Pi_YD\Pi_Y=\int_HU_h\Pi_YD\Pi_YU_h^\dagger\,dh=\Pi_Y\big(\int_HU_hDU_h^\dagger dh\big)\Pi_Y=0$ by \eqref{eq:W-twirl}.

\emph{(c)} Suppose $P_X=P_H$. By (a) some $Y\in\mathcal Y_H$ has $M:=Z_Y-D\succeq0$. Since $Z_Y\Pi_Y=0$ and the hypothesis of (b) holds for this $Y$,
$\Pi_YM\Pi_Y=-\Pi_YD\Pi_Y=0$. By Lemma~\ref{lem:zeroblock}, $0=(1-\Pi_Y)M\Pi_Y=(1-\Pi_Y)Z_Y\Pi_Y-(1-\Pi_Y)D\Pi_Y=-(1-\Pi_Y)D\Pi_Y$,
contradicting the hypothesis $(1-\Pi_Y)D\Pi_Y\ne0$. Hence $P_X\ne P_H$, and by (a) $P_X>P_H$.

\emph{(d)} (Here $Y\in\mathcal Y_H$, as implicit in the notation $Z_Y,\Pi_Y$.) With respect to $\mathcal K=\ker Z_Y\oplus(\ker Z_Y)^\perp$, $Z_Y$ is block
diagonal with blocks $0$ and $Z_Y^\perp:=(1-\Pi_Y)Z_Y(1-\Pi_Y)$; the hypotheses give $\Pi_YD\Pi_Y=0$ and $(1-\Pi_Y)D\Pi_Y=0$, hence also
$\Pi_YD(1-\Pi_Y)=((1-\Pi_Y)D\Pi_Y)^\dagger=0$. So $Z_Y-D$ is block diagonal with blocks $0$ and $Z_Y^\perp-(1-\Pi_Y)D(1-\Pi_Y)$. On $(\ker Z_Y)^\perp$,
$Z_Y^\perp\succeq\lambda^+_{\min}(Z_Y)$ and $(1-\Pi_Y)D(1-\Pi_Y)\preceq\lVert(1-\Pi_Y)D(1-\Pi_Y)\rVert\le\lambda^+_{\min}(Z_Y)$, so the second block is PSD.
Thus $Z_Y-D\succeq0$, and (a) gives $P_X=P_H$.

\emph{(e)} Let $M_n(X)=\E_{x\in X}(xx^\dagger)^{\otimes n}$ and $M_n(H)=\int_H(hf\,(hf)^\dagger)^{\otimes n}dh$ in $\mathrm{End}(V^{\otimes n})$, and
$m=M_n(X)-M_n(H)$. Let $\tau_T$ be the transpose on the first $T$ tensor factors, $\tau_T(A_1\otimes A_2)=A_1^{\rm t}\otimes A_2$ for
$A_1\in\mathrm{End}(V^{\otimes T})$, $A_2\in\mathrm{End}(V^{\otimes(T+k)})$, extended linearly; we read its output as an operator on
$\bar V^{\otimes T}\otimes V^{\otimes(T+k)}$ (same underlying space and basis). Let $P_{\mathcal K}$ be the orthogonal projection onto $\mathcal K$ and
$\mathcal L(A)=P_{\mathcal K}\,\tau_T(A)\,P_{\mathcal K}$, viewed in $\mathrm{End}(\mathcal K)$.

\emph{$D=\mathcal L(m)$.} $(xx^\dagger)^{\rm t}=\bar x\bar x^\dagger$, so $\tau_T((xx^\dagger)^{\otimes n})=u_xu_x^\dagger$, which already lives on $\mathcal K$; by linearity
(and continuity, for the Haar integral) $\mathcal L(M_n(X))=R_X$, $\mathcal L(M_n(H))=R_H$, and $D=\mathcal L(m)$.

\emph{Equivariance.} For product operators, $(hA_1h^\dagger)^{\rm t}=\bar h\,A_1^{\rm t}\,h^{\rm t}=\bar hA_1^{\rm t}\bar h^\dagger$ (as $h^{\rm t}=\bar h^\dagger$), so
$\tau_T(\mathrm{Ad}_h A)=(\bar h^{\otimes T}\otimes h^{\otimes(T+k)})\,\tau_T(A)\,(\bar h^{\otimes T}\otimes h^{\otimes(T+k)})^\dagger$ for all $A$ by linearity. $\mathcal K$ is
invariant under the unitary $\bar h^{\otimes T}\otimes h^{\otimes(T+k)}$, so $P_{\mathcal K}$ commutes with it, and $\mathcal L(\mathrm{Ad}_hA)=U_h\mathcal L(A)U_h^\dagger$.
For $Y\in\mathcal Y_H$ define $\Phi_Y:\mathrm{End}(V^{\otimes n})\to\mathrm{Hom}(\ker Z_Y,(\ker Z_Y)^\perp)$, $\Phi_Y(A)=(1-\Pi_Y)\mathcal L(A)|_{\ker Z_Y}$. Since $\Pi_Y$
commutes with $U_h$,
\[
\Phi_Y(\mathrm{Ad}_hA)=(1-\Pi_Y)U_h\mathcal L(A)U_h^\dagger\Pi_Y=U_h|_{(\ker Z_Y)^\perp}\,\Phi_Y(A)\,U_h|_{\ker Z_Y}^\dagger ,
\]
i.e.\ $\Phi_Y$ is $H$-equivariant for the actions (A3), and $(1-\Pi_Y)D\Pi_Y=\Phi_Y(m)$ (extended by $0$ on $(\ker Z_Y)^\perp$).

\emph{Isotypic decomposition with multiplicities.} $H$ is compact, so the finite-dimensional modules $\mathcal E=\mathrm{End}(V^{\otimes n})$ and
$\mathcal H_Y=\mathrm{Hom}(\ker Z_Y,(\ker Z_Y)^\perp)\cong(\ker Z_Y)^\perp\otimes(\ker Z_Y)^\ast$ decompose into isotypic components,
$\mathcal E=\bigoplus_\lambda\mathcal E_\lambda$ with $\mathcal E_\lambda\cong V_\lambda\otimes\C^{\mu_\lambda}$, and similarly $\mathcal H_Y=\bigoplus_\lambda(\mathcal H_Y)_\lambda$; the
isotypic projections $P_\lambda$ are $H$-equivariant (they are $\int_H\overline{\chi_\lambda(h)}\dim V_\lambda\,\mathrm{Ad}_h\,dh$). By Schur's lemma, an equivariant
map sends $\mathcal E_\lambda$ into $(\mathcal H_Y)_\lambda$, whatever the multiplicities: restricted to any irreducible summand $\cong V_\lambda$ of $\mathcal E_\lambda$, its
components into the irreducible summands of $\mathcal H_Y$ vanish unless they are $\cong V_\lambda$. Hence $\Phi_Y(m_\lambda)=0$ unless
$V_\lambda\hookrightarrow\mathcal H_Y$, where $m_\lambda=P_\lambda m$. The dual $(\ker Z_Y)^\ast$ enters only through this Hom space; for self-dual $H$-modules
(e.g.\ all $SU(2)$- and $SO(4)$-modules) $(\ker Z_Y)^\ast\cong\ker Z_Y$.

\emph{Trivial component.} $P_{\mathbf 1}=\int_H\mathrm{Ad}_h\,dh$, and $P_{\mathbf 1}M_n(X)=M_n(H)$ by the argument of \eqref{eq:W-twirl} (each $x\in X$ lies in
$H\cdot f$), while $P_{\mathbf 1}M_n(H)=M_n(H)$; so $m_{\mathbf 1}=0$.

\emph{$G$-fixed vectors.} $M_n(X)$ is $\mathrm{Ad}_G$-invariant and $M_n(H)$ is $\mathrm{Ad}_H$-invariant, so $m$ is $G$-fixed; $P_\lambda$ commutes with
$\mathrm{Ad}_g$, so $m_\lambda$ is a $G$-fixed vector of $\mathcal E_\lambda\cong V_\lambda\otimes\C^{\mu_\lambda}$, whose $G$-fixed space is $V_\lambda^G\otimes\C^{\mu_\lambda}$.
Hence $m_\lambda\ne0$ requires $V_\lambda^G\ne0$.

Therefore $(1-\Pi_Y)D\Pi_Y=\sum_\lambda\Phi_Y(m_\lambda)$ only receives contributions from $\lambda\ne\mathbf 1$ with $V_\lambda^G\ne0$ and
$V_\lambda\hookrightarrow\mathrm{Hom}(\ker Z_Y,(\ker Z_Y)^\perp)$. Since $m=M_{2T+k}(X)-M_{2T+k}(H)$, this is (e).
\end{proof}

\begin{corollary}[W-SU(2)]\label{cor:WSU2}
Let $H=SU(2)$ act on $V=V_j$ (spin $j$), $G\le H$ closed, and let $Y\in\mathcal Y_H$ with $\ker Z_Y$ spin-$s$ isotypic. Then
$(1-\Pi_Y)D\Pi_Y$ only receives contributions from spins $L$ with $1\le L\le\min\{2j(2T+k),\ s+j(2T+k)\}$ and $V_L^G\ne0$.
\end{corollary}
\begin{proof}
$V_j^{\otimes n}$ contains only spins $\le jn$, so $\mathrm{End}(V_j^{\otimes n})\cong V_j^{\otimes n}\otimes(V_j^{\otimes n})^*$ contains only spins
$L\le2jn$ (every $SU(2)$-irreducible is self-dual). $\bar V_j\cong V_j^\ast\cong V_j$, so $\mathcal K\subset\bar V_j^{\otimes T}\otimes V_j^{\otimes(T+k)}$ contains only
spins $p\le jn$, and so does $(\ker Z_Y)^\perp$. Write $\ker Z_Y\cong V_s\otimes\C^\mu$ and $(\ker Z_Y)^\perp\cong\bigoplus_pV_p\otimes\C^{\nu_p}$; then
$\mathrm{Hom}(\ker Z_Y,(\ker Z_Y)^\perp)\cong\bigoplus_p(V_p\otimes V_s^\ast)\otimes\C^{\mu\nu_p}$, and multiplicities create no new irreducibles. By the
triangle rule, $V_L\hookrightarrow V_p\otimes V_s^\ast\cong V_p\otimes V_s$ requires $L\le s+p\le s+jn$. $L=0$
is excluded by (e). Apply Theorem~\ref{thm:window}(e). (No use of (C1).)
\end{proof}

\begin{corollary}[W-O(4)]\label{cor:WO4}
Let $H=O(4)$ act on $V=\C^4$ by its defining real orthogonal representation, let $f\in\mathbb{R}^4$ be a real unit vector and $G\le O(4)$ finite,
and suppose $X=G\cdot f\subset S^3$ is a spherical $t$-design. Let $Y\in\mathcal Y_H$ and suppose $\ker Z_Y$ is isotypic, as an
$SO(4)$-module, of type $(a,b)$ under $SO(4)\cong(SU(2)\times SU(2))/\{\pm1\}$. Then $(1-\Pi_Y)D\Pi_Y=\sum_\ell(1-\Pi_Y)\mathcal L(m_\ell)\Pi_Y$,
where the sum runs over even $\ell$ with
\[
t<\ell\le2(2T+k),\qquad\ell\le2\min(a,b)+2T+k,
\]
and $m_\ell$ is the degree-$\ell$ spherical-harmonic part of $m$ defined in the proof. In particular, if no even $\ell$ satisfies both
bounds, then $(1-\Pi_Y)D\Pi_Y=0$.
\end{corollary}
\begin{proof}
\emph{Step 1 (harmonic parts of $m$).} For real $x$, $xx^\dagger=xx^{\rm t}$ and $\phi(x):=(xx^{\rm t})^{\otimes n}$ has entries that are
homogeneous polynomials of degree $2n$ in $x$. On $S^3$ each entry decomposes into spherical harmonics of degrees $\ell\le2n$, $\ell$ even
(the entries are even functions); write $\phi=\sum_\ell\phi_\ell$ entrywise. Since $O(4)$ commutes with the harmonic projections and
$\phi(hx)=h^{\otimes n}\phi(x)h^{{\rm t}\otimes n}$, also $\phi_\ell(hx)=h^{\otimes n}\phi_\ell(x)h^{{\rm t}\otimes n}$. $O(4)$ acts
transitively on $S^3$, so $M_n(H)=\int_{S^3}\phi\,d\sigma$, and $m=\sum_\ell m_\ell$ with
$m_\ell:=\E_X\phi_\ell-\int_{S^3}\phi_\ell\,d\sigma$. For $\ell\ge1$, $\int\phi_\ell\,d\sigma=0$; for $1\le\ell\le t$, $\E_X\phi_\ell=0$
by the design property; and $m_0=0$. So $m=\sum_{t<\ell\le2n,\ \ell\ \rm even}m_\ell$ with $m_\ell=\E_X\phi_\ell$. Let
$\mathcal H_\ell$ be the space of degree-$\ell$ harmonics, $Z^{(\ell)}_y\in\mathcal H_\ell$ the reproducing kernel at $y$, and
$\Psi_\ell(p)=\int_{S^3}p(x)\phi(x)\,d\sigma(x)$ for $p\in\mathcal H_\ell$. $\Psi_\ell$ is $SO(4)$-equivariant and
$\phi_\ell(y)=\Psi_\ell(Z^{(\ell)}_y)$, so $m_\ell=\Psi_\ell(\E_XZ^{(\ell)}_y)\in\Psi_\ell(\mathcal H_\ell)$. Since
$\mathcal H_\ell\cong(\ell/2,\ell/2)$ is an irreducible $SO(4)$-module, $m_\ell$ lies in the $(\ell/2,\ell/2)$-isotypic component.

\emph{Step 2 (types in $\mathcal K$).} $V\otimes_{\mathbb R}\C$ is $(\tfrac12,\tfrac12)$, and $\Sym^i(\C^4)=\bigoplus_{i'\le i,\ i'\equiv i\ (2)}\mathcal H_{i'}\cong
\bigoplus(i'/2,i'/2)$. $\bar V\cong V$, so by Clebsch--Gordan in each $SU(2)$ factor $\mathcal K=\Sym^T(\bar V)\otimes\Sym^{T+k}(V)$ contains only
types $(p,p')$ with $p,p'\le(T+T+k)/2=n/2$; so does $(\ker Z_Y)^\perp\subset\mathcal K$.

\emph{Step 3 (the Hom space).} The map $\Phi_Y$ of the proof of (e) is $O(4)$-equivariant, hence $SO(4)$-equivariant, so by the Schur argument
of (e) (applied to $SO(4)$, with multiplicities) it kills $m_\ell$ unless $(\ell/2,\ell/2)\hookrightarrow\mathrm{Hom}(\ker Z_Y,(\ker Z_Y)^\perp)\cong(\ker Z_Y)^*\otimes(\ker Z_Y)^\perp$.
$SU(2)$-irreducibles are self-dual, so $(a,b)^*\cong(a,b)$, and $(\ell/2,\ell/2)\hookrightarrow(a,b)\otimes(p,p')$ requires
$\ell/2\le a+p$ and $\ell/2\le b+p'$. With $p,p'\le n/2$ this gives $\ell\le2\min(a,b)+n$.
\end{proof}

\begin{remark}
Corollaries~\ref{cor:WSU2} and~\ref{cor:WO4} only restrict where $(1-\Pi_Y)D\Pi_Y$ can be nonzero. To conclude $P_X=P_H$ one still needs
(b) (or $\Pi_YD\Pi_Y=0$ directly) and the norm condition of (d). Applying them to a concrete orbit (e.g.\ the $O(4)$ row of Table~\ref{tab:nonab}) requires $Y$ and the type of $\ker Z_Y$, which
were obtained numerically [Num]. In the five $O(4)$ tests of the research phase the computed window degrees $4,6,8,8,8$ equal
$2\min(a,b)+2T+k$ and lie below the first surviving harmonic ($6$ for the $F_4$ orbit, $12$ for the $H_4$ orbit).
\end{remark}

\begin{table}[t]\centering
\caption{Finite-orbit tests at matched $(T+k)$-th moment (genuine cases: Lemma~\ref{lem:B} is not tight).}\label{tab:nonab}
\begin{tabular}{llccl}\toprule
$H$ / fiducial & $G$-orbit & $(T,k)$ & $P_X-P_H$ & status\\\midrule
$SO(3)$, real $\ket{1,0}$ & 6 icosahedral axes & (1,1) & $<10^{-13}$ & pass\\
$SO(3)$, real $\ket{1,0}$ & 30-line icosahedral orbit & (2,1) & $<10^{-12}$ & pass\\
$SO(3)$, real $\ket{1,0}$ & 30-line icosahedral orbit & (1,3) & $<10^{-11}$ & pass\\
$SU(2)$, spin $3/2$ GHZ-like & icosahedral (tuned) & (1,1) & $<10^{-12}$ & pass ($P=31/70$)\\
$O(4)$, real & 24-cell & (1,1) & $<10^{-8}$ & pass\\
$SU(2)$, spin 1 (via $SO(3)$) & octahedral (tuned) & (1,1) & $\ge6.437\times10^{-4}$ & fail (certified)\\
$SU(2)$, spin $5/2$ & binary icosahedral (tuned) & (1,1) & $1.1$--$8.8\times10^{-4}$ & fail (predicted by Thm.~\ref{thm:window})\\
$SU(3)$, $\Sym^2\C^3$ & Hessian (tuned) & (1,1) & $\ge8.0399\times10^{-4}$ & fail (certified)\\\bottomrule
\end{tabular}
\end{table}

\section{Related work}\label{sec:related}

\paragraph{Optimal cloning for continuous families.}
Werner~\cite{Werner1998} constructed the unique optimal $N\to M$ cloner for pure $d$-level states under global fidelity.
For equatorial (phase) states, Bru\ss{} \emph{et al.}~\cite{Bruss2000} bounded the $N\to M$ fidelity via phase estimation and gave the optimal $1\to2$ qubit machine;
Fan \emph{et al.}~\cite{Fan2003} obtained the optimal $1\to2$ phase-covariant cloner for qudits under single-qudit fidelity;
and Buscemi, D'Ariano and Macchiavello~\cite{Buscemi2005} derived the optimal $N\to M$ phase-covariant cloner for $M=kd+N$ (their $k$ is an integer multiplier, not our number of extra copies), optimal for both global and single-qudit fidelity.
Kim and Chitambar~\cite{KimChitambar2022} computed the optimal global (process) fidelity of phase-covariant cloning by semidefinite programming, with closed form $(2d-1)/d^{2}$ for $1\to2$, and asked for a $1\to M$ formula.
That formula, and in fact the optimal $N\to N+K$ phase-state cloner for all $d,N,K$, is now given by Herasymenko, Lin, Verduyn Lunel and Grinko~\cite[Thm.~II.3 and Thm.~X.1, Eq.~(16)]{Herasymenko2026}.
In particular the $T=1$ torus values, $(k+1)!(2d-k)/(2d^{k+1})$ for $k\le d$, are the $N=1$ case of their Eq.~(16); we use their result for the torus value and for the singleton case of Proposition~\ref{prop:reduction}, and claim no novelty for it.
Asymptotic and asymmetric variants are treated by Chiribella and Yang~\cite{ChiribellaYang2014}, who relate optimal asymptotic cloners to state estimation for, among others, phase- and multiphase-covariant families, and by Kay~\cite{Kay2016} for universal cloning.
\emph{Difference.} All of these works optimise over the full continuous orbit (the torus of phases, or the unitary group).
We ask a different question: which \emph{finite} phase codebooks $\{\psi_x\}_{x\in\mathbb{Z}_q}$ with key set $K\subset\mathbb{Z}_q$ attain the torus value $P_{\mathrm{torus}}$ exactly.

\paragraph{Finite sets that reproduce a continuous optimum.}
The closest precedent is Jing \emph{et al.}~\cite{Jing2012}: three equally spaced equatorial qubit states determine the phase-covariant cloner, and four tetrahedral states determine the universal one.
Their statement is for qubits and single-copy fidelity.
For $d=2$, $T=k=1$, our sufficiency result gives the same conclusion under global fidelity, since $\{0,1\}\subset\mathbb{Z}_3$ is a $B_2$ set.
Our necessity statement (Theorem~\ref{thm:sidon}: Sidon is required at $(T,k)=(1,1)$) holds \emph{for global fidelity only} and does not transfer to single-copy fidelity, where Jing \emph{et al.} report that three states suffice for all $1\to n$. For larger $(T,k)$ we prove sufficiency of $B_{T+k}$ only; necessity is open.
\'Cwikli\'nski and Studzi\'nski~\cite{CwiklinskiStudzinski2026} use moment matching in the opposite direction: when the second moment of an orbit equals the Haar moment, the orbit is a continuous 2-design and local cloning reduces to universal cloning.
Herasymenko \emph{et al.}~\cite{Herasymenko2026} show that stabilizer states are not a resolving orbit by a 4-copy frame-potential computation; Gross, Nezami and Walter~\cite{GrossNezamiWalter2021} supply the Clifford-group Schur--Weyl duality that controls such stabilizer moments.
\emph{Difference.} We give (i) an exact finite-codebook criterion on the torus: $B_{T+k}$ index sets attain $P_{\mathrm{torus}}$, proved for $d=2$ and for $T\le3$ (all $d,k$), with an if-and-only-if (Sidon) characterisation at $(T,k)=(1,1)$; and (ii) a window theorem (Theorem~\ref{thm:window}) for finite orbits of compact groups, with two computer-assisted counterexamples (SU(3) on $\mathrm{Sym}^2\mathbb{C}^3$ with the Hessian group, and spin-1 with the octahedral group) showing that the finite-subgroup analogue fails beyond tori.
We found no prior finite-orbit failure theorem or window criterion.

\paragraph{Quantum fingerprinting and hashing.}
Buhrman, Cleve, Watrous and de~Wolf~\cite{Buhrman2001} introduced quantum fingerprints that are exponentially shorter than the strings they encode; Gottesman and Chuang~\cite{GottesmanChuang2001} built unconditionally secure quantum digital signatures from quantum public keys, whose security requires that only a limited number of copies circulate.
Gavinsky and Ito~\cite{GavinskyIto2013} constructed \emph{hiding} quantum fingerprints that reveal $o(\min\{n,\log(1/\varepsilon)\})$ bits.
Ablayev and Vasiliev~\cite{AblayevVasiliev2014} introduced cryptographic quantum hashing with a Fourier-controlled key set and a trigonometric (qubit-rotation) encoding; Vasiliev~\cite{Vasiliev2016} formulated the character-based phase encoding $\frac1{\sqrt d}\sum_j e^{2\pi i s_jx/q}\lvert j\rangle$ for finite abelian groups, which is exactly our state family, with an $\varepsilon$-biased key set $S\subset\mathbb{Z}_q$ giving collision resistance (see also the review~\cite{Ablayev2025}); Ziiatdinov~\cite{Ziiatdinov2016} extends the construction to arbitrary groups.
Hatanaka \emph{et al.}~\cite{Hatanaka2024} propose a photonic hash based on Gaussian boson sampling.
Adaptive attacks on Wiesner money~\cite{Nagaj2016} illustrate how copying and verification interact in practice.
\emph{Difference.} These works use Fourier bias, i.e.\ near-orthogonality, for collision resistance.
None of them analyses optimal \emph{cloning} of the key-set states, and none uses a Sidon or $B_h$ condition.
Our results give the corresponding copy-resistance statement: for $B_{T+k}$ key sets, an adversary holding $T$ copies does no better at producing $T+k$ copies (in global fidelity) than against the full phase torus.

\paragraph{Sidon sets, difference sets and their uses in quantum information.}
Sidon ($B_2$) sets go back to Erd\H{o}s and Tur\'an~\cite{ErdosTuran1941}.
Singer~\cite{Singer1938} gave perfect difference sets modulo $q^2+q+1$, and Bose and Chowla~\cite{BoseChowla1962} extended these results to $B_h$ sequences.
Katz's character-sum estimate~\cite{Katz1989} bounds $\lvert\sum_{t\in F}\chi(t-x)\rvert\le(n-1)\sqrt{\#F}$ for $x$ of degree $n$ over $F$; it controls the Fourier bias of Bose--Chowla-type sets, which we use in the Discussion.
O'Bryant's survey~\cite{OBryant2004} is the standard bibliography.
In quantum information, Bertoni, Mereghetti and Palano~\cite{BertoniMereghettiPalano2003} used Golomb rulers and difference sets to minimise the size of one-way quantum finite automata, and Aydin, Albert and Barg~\cite{AydinAlbertBarg2026} built $\ell_1$ codes, and hence spin, bosonic and permutation-invariant quantum codes, from $t$-Sidon sets and Tverberg partitions.
Vitelli \emph{et al.}~\cite{Vitelli2013} realised quantum state fusion of two photonic qubits into one ququart.
\emph{Difference.} In \cite{BertoniMereghettiPalano2003,AydinAlbertBarg2026} the additive structure serves automaton size or code distance.
Here it is the exact condition for finite codebooks to attain optimal cloning.
Our identity between the optimal fusion fidelity of phase states and the additive energy of the key set (Theorem~\ref{thm:fusion}) has no counterpart in~\cite{Vitelli2013}, which is a physical protocol.
The multiset Johnson-type packing bound used in the proofs of Theorem~\ref{thm:packing}(c,d) is a packing statement; it is not the Erd\H{o}s--Ko--Rado bound for intersecting multiset families of Meagher and Purdy~\cite{MeagherPurdy2011}, which goes in the opposite direction.

\section{Discussion}
For quantum hashing and copy-bounded quantum tokens, our results replace the bias heuristic by an exact criterion: at $(T,k)=(1,1)$ a phase codebook is as hard to counterfeit as the full torus exactly when its key set is Sidon, and for larger
$(T,k)$ the absence of additive coincidences of order $T+k$ ($B_{T+k}$) suffices (proved for $d=2$ and for $T\le3$; conjectured in
general); whether it is also necessary is open. Bose--Chowla sets~\cite{BoseChowla1962} are $B_h$ with small bias~\cite{Katz1989}, but they are not efficiently
indexable on qubit registers, and the scheme remains malleable and vulnerable to verification-oracle attacks; these results are
information-theoretic statements about copy-bounded security, not a deployable primitive.
\paragraph{Open problems.} (1) Conjecture~\ref{conj:TT} for $T\ge4$; the certified margin $1-\max\rho\approx2.9/T^2$ at $d=3$ shows that a proof
must be sharp in $T$. (2) Exact conditions for $T\ge2$: Theorem~\ref{thm:t1iff} settles $T=1$ (and shows that $B_{T+k}$ is not necessary); is there an analogous
characterisation for $T\ge2$, and does Conjecture~\ref{conj:t1H} hold for all $(d,k)$? (3) A necessary-and-sufficient form of
Theorem~\ref{thm:window}, i.e.\ a characterisation of the finite orbits whose cloning value equals that of the continuous orbit. (4) Codebooks
over other finite abelian groups (e.g.\ $\Z_q^r$), and cloning figures of merit other than global fidelity.

\section*{Use of AI tools}
This work made substantial use of four AI systems.
\begin{itemize}
\item \textbf{Grok Bot} (xAI's agent platform, launched on the Grok~4.6 model; used October 2026) was used throughout the research phase. It proposed and drafted proofs
(the reduction to the packing inequality, the packing inequality for $d=2$ and $T\le3$, the fusion theorems, the $T=1$ characterisation and the window theorem),
wrote exploratory numerical code and the certificate generators, wrote the standalone exact, interval and ball-arithmetic checkers of
Appendix~\ref{app:ca}, and ran literature and prior-art searches.
\item \textbf{Claude} (Anthropic; Claude Opus 5.5) co-developed research directions and conjectures, independently re-implemented
and audited every computational claim in a separate code base (QH-Lab~\cite{QHLab}), re-ran every certificate checker, verified
every bibliographic entry against Crossref and arXiv, and drafted and edited the manuscript.
\item \textbf{ChatGPT} (OpenAI; GPT-6 Luna, the ChatGPT Free model since 8 October 2026; used 16 October 2026) served as an independent referee of the draft, re-ran the standard-library certificate checkers, and gave an independent proof of
Theorem~\ref{thm:fusionm}, and \textbf{Perplexity}
(Perplexity Pro with the GPT-6.1 Sol model) as an independent prior-art and citation-fidelity checker. Every issue they raised was adjudicated by
recomputation or by a written proof before any change was made.
\end{itemize}
The outputs were verified as follows. (1) Every computer-assisted step is checked by a standalone program that uses only exact
rational, interval or ball arithmetic and no optimisation solver; floating-point and solver outputs serve only as certificate data.
(2) The complete certificate package was re-run independently: the archive hash and all 32 internal hashes matched, and every output
was identical to the shipped expected output. (3) The values in Tables~\ref{tab:torus}--\ref{tab:nonab} and in the negative-results tables, the worked examples, the $m$-fold fusion values
and the $T=1$ characterisation (on random key sets) were reproduced by an independent implementation (QH-Lab), together with floating-point spot checks of the torus theorem at $T=2$ ($d\le20$) and $T=3$ ($d\le8$), which agree with the
torus value to within $10^{-9}$. (4) The pen-and-paper arguments were reviewed line by line by Claude against the research record, and by two independent ChatGPT
referee passes on successive drafts; the author read the complete manuscript.
(5) Every reference was checked against its arXiv or DOI landing page.
The author takes full responsibility for the content of this paper, including all text and results produced with AI assistance.

\section*{Data and code availability}
All code needed to verify the computer-assisted steps is provided as arXiv ancillary files and in the Zenodo record that accompanies
this paper~\cite{QHLab} (code under the MIT licence). The archive \texttt{ca\_appendix.zip} has SHA-256
\texttt{e77e131bdd7becb0822ed582ccfc0a9e43b4b9183aeab37e65bec3e64d888b30}; the separate coverage audit \texttt{coverage\_check.py} has SHA-256
\texttt{c42d11f652ed0bdd7daec55dd4bb9940319a8ff138c7f7d59956eacfa9ef0995}. The values in the negative-results tables are reproduced by
QH-Lab's \texttt{tools/neg\_table\_check.py}. Verification requires Python $\ge3.10$,
\texttt{mpmath} and \texttt{python-flint} 0.9.0. No other data were generated or analysed.

\section*{Author contributions}
M.~Weldon is the sole author. He conceived the project, directed the research and its multi-system verification workflow, read and
approved the final manuscript, and takes responsibility for its content, with AI assistance as described above.

\section*{Competing interests}
The author declares no competing interests. This work received no external funding.

\bibliographystyle{unsrtnat}
\bibliography{refs}

\appendix
\providecommand{\Mb}{\mathrm{Mb}}
\providecommand{\ov}{\mathrm{ov}}
\providecommand{\checker}[1]{\textsf{[CA:~\nolinkurl{#1}]}}
\section{Proof of the packing inequality (Theorem~\ref{thm:packing})}\label{app:proofs}

\paragraph{Notation and scope.} Throughout, $N=T+k$. A \emph{shift} is $\Delta\in\Z^d$ with $\sum_i\Delta_i=k$; $\Delta=\Delta^+-\Delta^-$ with
$\Delta^\pm\ge0$ of disjoint supports, and $n(\Delta)=|\Delta^-|$. The \emph{support} of $\Delta$ is $\mathcal D_\Delta=\{b:|b|=T,\ b\ge\Delta^-\}$
(this is $\mathcal D_{\Delta}$ of~\cite[Eq.~(111)]{Herasymenko2026}) and its \emph{outputs} are $b+\Delta$, $b\in\mathcal D_\Delta$.
A \emph{family} is a finite set $F$ of shifts with $n(\Delta)\le T$ for all $\Delta\in F$ that pairwise satisfy
\[
\text{(ii)}\qquad\lVert\Delta-\Delta'\rVert_1\ge2(T+k+1).
\]
For $x,y\in\Z_{\ge0}^d$: $|x|=\sum_ix_i$, $\ov(x,y)=\sum_i\min(x_i,y_i)$, $\mult(x)=|x|!/\prod_ix_i!$. $\Mb(M,\ell)$ is the largest multinomial
$\mult(x)$ over $x$ with $|x|=M$ and at most $\ell$ nonzero coordinates (attained by the balanced vector), $M^\ast=\Mb(N,d)$, and
$R_m=m!/m^m$. We write $c_b=\mult(b)$ for $|b|=T$, so $a_b=\sqrt{c_b/Z_b}$. The tag \checker{path} marks a step certified by the named
standalone checker of Appendix~\ref{app:ca} (paths relative to \texttt{ca\_appendix/}); every other step is proved in the text.

\begin{remark}[Scope of ``family'']\label{rem:famscope}
A shift with $n(\Delta)>T$ has $\mathcal D_\Delta=\emptyset$ and $\Lambda(\Delta)=0$, and it can be at distance $\ge2(T+k+1)$ from a balanced $g$;
e.g.\ $d=3$, $T=k=1$, $g=(1,0,0)$, $\Delta'=(-1,-1,3)$, $\lVert g-\Delta'\rVert_1=6$. Then $\rho(\{g,\Delta'\})=1$. So parts (b)--(d) of
Theorem~\ref{thm:packing} (and the strict inequality of Conjecture~\ref{conj:TT}) hold only for families with $n(\Delta)\le T$, which is how they
are proved here; for Proposition~\ref{prop:reduction} the restriction is harmless, since members with $\Lambda=0$ do not change $\rho$.
Part (a) is proved below without the restriction.
\end{remark}

\subsection{Common lemmas}

\begin{lemma}[Normalisation, Lemma N]\label{lem:N}
Let $\mathrm{avg}_b=|C|^{-1}\sum_{g\in C}\mult(b+g)$, $\theta_b=\sqrt{c_b/\mathrm{avg}_b}$ and $W=\sum_{|b|=T}\sqrt{c_b\,\mathrm{avg}_b}$. Then for
every finite set $F$ of shifts
\[
\rho(F)=\frac1W\sum_{\Delta\in F}\ \sum_{b\in\mathcal D_\Delta}\theta_b\,\mult(b+\Delta).
\]
\end{lemma}
\begin{proof}
$Z_b=|C|\,\mathrm{avg}_b$, so $a_b=\theta_b/\sqrt{|C|}$. The $a_b$ are $S_d$-invariant and $C$ is one $S_d$-orbit, so $\Lambda(g)=\Lambda_0$ for
all $g\in C$; averaging over $C$, $\Lambda_0=\sum_ba_b\,\mathrm{avg}_b=|C|^{-1/2}\sum_b\theta_b\mathrm{avg}_b=W/\sqrt{|C|}$. Divide
$\sum_F\Lambda(\Delta)=|C|^{-1/2}\sum_F\sum_{b}\theta_b\mult(b+\Delta)$ by $\Lambda_0$.
\end{proof}
Write $\kappa_b=\mathrm{avg}_b/M^\ast$. Then $\theta_b\mult(\cdot)/W=\sqrt{c_b/\kappa_b}\,(\mult(\cdot)/M^\ast)\big/\sum_b\sqrt{c_b\kappa_b}$, so all
bounds below may be computed in units of $M^\ast$ with $\theta_b=\sqrt{c_b/\kappa_b}$ and $W=\sum_b\sqrt{c_b\kappa_b}$.

\begin{lemma}[Structure, Lemma S]\label{lem:S}
For shifts $\Delta,\Delta'$ with $n=n(\Delta)$, $n'=n(\Delta')$: $\lVert\Delta\rVert_1=k+2n$, and
$\lVert\Delta-\Delta'\rVert_1=\lVert\Delta\rVert_1+\lVert\Delta'\rVert_1-2\ov(\Delta^+,\Delta'^+)-2\ov(\Delta^-,\Delta'^-)$. Hence
\[
\text{(ii)}\iff\ov(\Delta^+,\Delta'^+)+\ov(\Delta^-,\Delta'^-)\le n+n'-T-1 .
\]
Consequently, in a family: (S1) every pair has $n+n'\ge T+1$; (S2) at most one member has $n\le T/2$; (S3) if $|F|\ge2$, every member has
$n\ge1$; (S4) a member with $n=T$ is \emph{pinned}: $\mathcal D_\Delta=\{\Delta^-\}$; (S5) supports of distinct members are disjoint and
outputs of distinct members are distinct.
\end{lemma}
\begin{proof}
$|\Delta^+|-|\Delta^-|=k$ gives $\lVert\Delta\rVert_1=|\Delta^+|+|\Delta^-|=k+2n$. For integers $x=x^+-x^-$, $y=y^+-y^-$ with
$x^+x^-=y^+y^-=0$ one checks $|x-y|=|x|+|y|-2\min(x^+,y^+)-2\min(x^-,y^-)$ (if $x,y\ge0$ this is $|x-y|=x+y-2\min(x,y)$; if $x\ge0\ge y$
both minima vanish; the other cases are symmetric). Summing over coordinates gives the identity, and substituting into (ii) the
equivalence. (S1): the left side is $\ge0$. (S2): two members with $n,n'\le T/2$ would have $n+n'\le T$. (S3): by (S1) a partner of a
member with $n=0$ has $n'\ge T+1$, which is excluded in a family. (S4): $b\ge\Delta^-$ and $|b|=T=|\Delta^-|$ force $b=\Delta^-$. (S5): if
$b\in\mathcal D_\Delta\cap\mathcal D_{\Delta'}$, then $\beta=b+\Delta$, $\beta'=b+\Delta'$ are nonnegative with $|\beta|=|\beta'|=N$, so
$\lVert\Delta-\Delta'\rVert_1=\lVert\beta-\beta'\rVert_1\le2N$, contradicting (ii); if $b+\Delta=b'+\Delta'$, then
$\lVert\Delta-\Delta'\rVert_1=\lVert b'-b\rVert_1\le2T$, again a contradiction.
\end{proof}

\begin{lemma}[Multinomial moves, Lemma M]\label{lem:M}
Moving one unit from a coordinate of value $a$ to one of value $c$ multiplies $\mult$ by $a/(c+1)$. Consequently:
\begin{enumerate}
\item[(M1)] (Doubling) If $M>\ell$ and $\ell<d$, then $\Mb(M,\ell+1)\ge2\Mb(M,\ell)$. Hence $\Mb(N,\ell)\le M^\ast2^{\ell-d}$ for $\ell\le d\le N$,
and $\Mb(M,\ell)\le M!\,2^{\ell-M}$ for $1\le\ell\le M$.
\item[(M2)] (Filling) Let $m=\lfloor N/d\rfloor$ and let $|x|=N$ with some coordinate $x_z=t_0\le m$. Then
$\mult(x)/M^\ast\le m!/(t_0!\,m^{m-t_0})$. In particular $\Mb(N,\ell)/M^\ast\le R_m^{\,d-\ell}$ for $\ell\le d$.
\item[(M3)] ($f$-monotonicity) For $1\le s<M$ and $s<d$: $\Mb(M,s+1)/(s+1)\ge\frac{2s}{s+1}\Mb(M,s)/s\ge\Mb(M,s)/s$.
\end{enumerate}
\end{lemma}
\begin{proof}
The move rule is immediate from the factorials. (M1) A balanced vector with $\ell$ parts summing to $M>\ell$ has a part $a\ge2$; moving one
unit to an empty coordinate multiplies $\mult$ by $a\ge2$ and gives a vector with $\ell+1$ parts. Iterating from $\ell$ to $d$ (allowed
since $N\ge d$) gives the first bound. For the second, a vector with $\ell\le M$ positive parts $x_i$ has $\prod x_i!\ge\prod2^{x_i-1}=2^{M-\ell}$.
(M2) While $x_z=t<m$, the other $d-1$ coordinates sum to $N-t>N-m\ge(d-1)m$, so their maximum is $\ge m+1$; moving one unit from it to $z$
multiplies $\mult$ by $\ge(m+1)/(t+1)\ge m/(t+1)$. After $m-t_0$ moves the vector is still of size $N$, so $\mult\le M^\ast$, i.e.\
$\mult(x)\prod_{t=t_0}^{m-1}\frac m{t+1}\le M^\ast$, which is the claim. For the second statement, fill the $d-\ell$ empty coordinates one after
another ($t_0=0$ each time; already filled coordinates have value $m$ and are never the donor, since the donor has value $\ge m+1$). (M3)
By (M1), $\Mb(M,s+1)/(s+1)\ge2\Mb(M,s)/(s+1)$, and $2s\ge s+1$.
\end{proof}

\begin{lemma}[Johnson-type multiset packing, Lemma J]\label{lem:J}
Let $s_r(P)$ be the number of distinct $r$-sub-multisets of a multiset $P$ on $[d]$. If the members of a collection of multisets pairwise
satisfy $\ov(P,P')\le r-1$, then $\sum_Ps_r(P)\le\binom{d+r-1}{r}$. For $|P|=M$ with support size $\ell$ and $\ell_j=\#\{x:P_x\ge j\}$:
\begin{gather*}
s_2(P)=\tbinom\ell2+\ell_2\ge\tbinom\ell2+[M>\ell],\\
s_3(P)=\tbinom\ell3+\ell_2(\ell-1)+\ell_3\ge\tbinom\ell3+(\ell-1)[M>\ell]+[M\ge\ell+2]=:s_3^{\min}(\ell,M).
\end{gather*}
\end{lemma}
\begin{proof}
An $r$-multiset below both $P$ and $P'$ would give $\ov(P,P')\ge r$; so the sets of $r$-sub-multisets are pairwise disjoint subsets of the
$\binom{d+r-1}r$ $r$-multisets on $[d]$. The formulas count sub-multisets by shape ($\{x,y\}$, $\{x,x\}$; $\{x,y,z\}$, $\{x,x,y\}$,
$\{x,x,x\}$). $M>\ell$ forces $\ell_2\ge1$. If $M\ge\ell+2$, either some coordinate is $\ge3$ ($\ell_3\ge1$, $\ell_2\ge1$) or two coordinates
are $\ge2$ ($\ell_2\ge2$ and $2(\ell-1)\ge(\ell-1)+1$ since $\ell\ge2$).
\end{proof}
Lemma~\ref{lem:J} is the multiset analogue of the Johnson packing bound; it is not an Erd\H{o}s--Ko--Rado statement~\cite{MeagherPurdy2011}.

\subsection{Part (a): \texorpdfstring{$d=2$}{d=2}, all \texorpdfstring{$T,k$}{T,k}}\label{app:d2}
Throughout this subsection $d=2$ and $N=T+k$. Write $b=(b,T-b)$ with $b\in\{0,\dots,T\}$, $\Delta=(\delta,k-\delta)$ with $\delta\in\Z$, and
$c_m=\binom Nm$ for $0\le m\le N$, $c_m=0$ otherwise. The output $b+\Delta=(b+\delta,N-b-\delta)$ is nonnegative iff $0\le b+\delta\le N$, so
\begin{gather*}
\mathcal D_\Delta=\{b\in\Z:\ \max(0,-\delta)\le b\le\min(T,N-\delta)\}\quad\text{(an interval, possibly empty)},\\
\Lambda(\Delta)=\sum_{b=0}^Ta_b\,c_{b+\delta}.
\end{gather*}
The boundary cases $b=0$, $b=T$ and partial or empty supports are all covered by the convention $c_m=0$ for $m\notin[0,N]$.

\begin{lemma}[Dominated rearrangement]\label{lem:rearr}
Let $x_0\ge x_1\ge\dots\ge x_n\ge0$ and let $(y_m)_{m\in Y}$ be nonnegative numbers on a finite index set $Y$ with $|Y|\ge n+1$; let
$y_{(0)}\ge\dots\ge y_{(n)}$ be the $n+1$ largest of them (with multiplicity). If $U\subset\{0,\dots,n\}$ and $\sigma:U\to Y$ is injective, then
$\sum_{u\in U}x_u\,y_{\sigma(u)}\le\sum_{i=0}^nx_i\,y_{(i)}$.
\end{lemma}
\begin{proof}
Let $\gamma_0\ge\gamma_1\ge\dots$ be the values $y_{\sigma(u)}$, $u\in U$, sorted decreasingly. They sit at $|U|$ distinct positions of $Y$
(injectivity), so for each $i$ the $i+1$ values $\gamma_0,\dots,\gamma_i$ occupy $i+1$ distinct positions; at most $i$ positions of $Y$
carry a value $>y_{(i)}$, hence $\gamma_i\le y_{(i)}$. Likewise the sorted values $\xi_0\ge\xi_1\ge\dots$ of $(x_u)_{u\in U}$ satisfy
$\xi_i\le x_i$. By the rearrangement inequality and nonnegativity,
$\sum_{u\in U}x_uy_{\sigma(u)}\le\sum_i\xi_i\gamma_i\le\sum_{i<|U|}x_iy_{(i)}\le\sum_{i=0}^nx_iy_{(i)}$.
\end{proof}

\begin{proof}[Proof of Theorem~\ref{thm:packing}(a)]
Let $F$ be any finite set of shifts that pairwise satisfy (ii) (no condition on $n(\Delta)$ is needed when $d=2$; members with empty
support contribute $0$).

\emph{Step 1: separation gives a partial injection.} For $d=2$, $\lVert\Delta-\Delta'\rVert_1=|\delta-\delta'|+|(k-\delta)-(k-\delta')|=2|\delta-\delta'|$,
so (ii) reads $|\delta-\delta'|\ge N+1$. (i) Supports are disjoint: if $b\in\mathcal D_\Delta\cap\mathcal D_{\Delta'}$, then
$b+\delta,b+\delta'\in[0,N]$, so $|\delta-\delta'|\le N$. (ii) Outputs are distinct: if $b+\delta=b'+\delta'$ with $b,b'\in[0,T]$, then
$|\delta-\delta'|=|b-b'|\le T<N+1$. Let $U=\bigcup_{\Delta\in F}\mathcal D_\Delta\subset\{0,\dots,T\}$; by (i) each $b\in U$ lies in the support of
exactly one member $\Delta(b)=(\delta(b),k-\delta(b))$, and $\sigma(b):=b+\delta(b)\in[0,N]$ is injective on $U$ by (ii). Hence
\[
\sum_{\Delta\in F}\Lambda(\Delta)=\sum_{b\in U}a_b\,c_{\sigma(b)} .
\]

\emph{Step 2: the $T+1$ largest $c$-values, before any rearrangement.} Let $c_{(0)}\ge\dots\ge c_{(T)}$ be the $T+1$ largest entries of
$(c_m)_{m=0}^N$ ($N+1\ge T+1$ entries). Put $\delta_0=\lceil k/2\rceil$ and $g=(\delta_0,k-\delta_0)$; $g\in C$ in both parities ($g=(k/2,k/2)$ for $k$ even,
$g=(\frac{k+1}2,\frac{k-1}2)$ for $k$ odd), and $\mathcal D_g=\{0,\dots,T\}$ because $g\ge0$. The window $\{\delta_0,\dots,\delta_0+T\}$ of row $N$
is $\{\lceil k/2\rceil,\dots,N-\lfloor k/2\rfloor\}$. Every window position $i$ has $i\ge\lceil k/2\rceil\ge\lfloor k/2\rfloor$ and $N-i\ge\lfloor k/2\rfloor$,
so $\min(i,N-i)\ge\lfloor k/2\rfloor$. Every omitted position $j$ has either $j\le\lceil k/2\rceil-1\le\lfloor k/2\rfloor$ or $N-j\le\lfloor k/2\rfloor-1$, so
$\min(j,N-j)\le\lfloor k/2\rfloor$. Since $\binom Nm$ is a nondecreasing function of $\min(m,N-m)$, every window value is $\ge$ every omitted value;
hence the multiset $\{c_{b+\delta_0}\}_{b=0}^T$ is exactly $\{c_{(0)},\dots,c_{(T)}\}$.

\emph{Step 3: upper bound.} By Lemma~\ref{lem:rearr} with $x=(a_b)$ sorted decreasingly, $Y=\{0,\dots,N\}$, $y=c$ and the injection $\sigma$ of Step 1,
$\sum_F\Lambda\le\Lambda^\ast:=\sum_{i=0}^Ta_{(i)}c_{(i)}$, where $a_{(0)}\ge\dots\ge a_{(T)}$ is the decreasing rearrangement of $(a_b)_{b=0}^T$.

\emph{Step 4: unimodality of $a_b$ (all $b$, including $b=0,T$).} For even $L\ge2$ and $0\le b\le T$ put $r_L(b)=\binom Tb/\binom{T+L}{b+L/2}$
(the lower index $b+L/2\in[L/2,T+L/2]$ is always admissible). Then $r_L(T-b)=r_L(b)$ and, for $0\le b\le T-1$,
\[
\frac{r_L(b+1)}{r_L(b)}=\frac{(T-b)(b+L/2+1)}{(b+1)(T+L/2-b)},\qquad (T-b)(b+\tfrac L2+1)-(b+1)(T+\tfrac L2-b)=\tfrac L2(T-1-2b),
\]
so $r_L$ is nondecreasing on $0\le b\le(T-1)/2$; with the symmetry, $r_L(b)$ is a nonincreasing function of $|b-T/2|$. If $k$ is even,
$C=\{(k/2,k/2)\}$, $Z_b=\binom N{b+k/2}$ and $a_b^2=r_k(b)$. If $k$ is odd, $C=\{(\frac{k+1}2,\frac{k-1}2),(\frac{k-1}2,\frac{k+1}2)\}$ and Pascal's rule
gives $Z_b=\binom N{b+(k+1)/2}+\binom N{b+(k-1)/2}=\binom{N+1}{b+(k+1)/2}$, so $a_b^2=r_{k+1}(b)$ (with $T+k+1=N+1$). In both cases $a_b$ is a
nonincreasing function of $|b-T/2|$.

\emph{Step 5: $\Lambda_0=\Lambda^\ast$.} $\Lambda_0=\Lambda(g)=\sum_{b=0}^Ta_bc_{b+\delta_0}$. By Step 2 the second factors are $c_{(0)},\dots,c_{(T)}$ in some
order, so it suffices that the pairing $b\mapsto(a_b,c_{b+\delta_0})$ is similarly ordered, i.e.\ that $a_b>a_{b'}$ implies $c_{b+\delta_0}\ge c_{b'+\delta_0}$
(then sorting by $a$ with ties broken by $c$ sorts both sequences, and $\Lambda_0=\Lambda^\ast$). Note $c_{b+\delta_0}$ is a nonincreasing function of
$|b+\delta_0-N/2|$, which is $|b-T/2|$ for $k$ even and $|b-\frac{T-1}2|$ for $k$ odd. If $a_b>a_{b'}$, then $|b-\frac T2|<|b'-\frac T2|$ (Step 4).
For $k$ even this is the claim. For $k$ odd, both numbers lie in $\frac T2+\Z$, so $|b-\frac T2|\le|b'-\frac T2|-1$, hence
$|b-\frac{T-1}2|\le|b-\frac T2|+\frac12\le|b'-\frac T2|-\frac12\le|b'-\frac{T-1}2|$, and $c_{b+\delta_0}\ge c_{b'+\delta_0}$.

Combining Steps 1--5, $\sum_F\Lambda\le\Lambda^\ast=\Lambda_0$, i.e.\ $\rho(F)\le1$, with equality for $F=\{g\}$.
\end{proof}

\subsection{Part (b): \texorpdfstring{$T=1$}{T=1}, all \texorpdfstring{$d\ge2$}{d>=2}, \texorpdfstring{$k\ge1$}{k>=1}}\label{app:T1}
Here $N=k+1$ and $|b|=1$, so $b=e_a$. We prove the slightly stronger bound $\rho(F)\le\frac{q+2}{2(d-1)(q+1)}$ for $k\ge d$, $q=\lfloor k/d\rfloor$,
which implies the stated $3/(4(d-1))$ because $q\ge1$.
\begin{proof}[Proof of Theorem~\ref{thm:packing}(b)]
Let $F$ be a family (Remark~\ref{rem:famscope}) with $|F|\ge2$.

\emph{Step 1: shape of the members.} Every member has $n(\Delta)\in\{0,1\}$. A member with $n=0$ and any partner with $n'\le1$ would need, by
Lemma~\ref{lem:S}, $\ov(\Delta^+,\Delta'^+)+\ov(\Delta^-,\Delta'^-)\le n+n'-T-1\le-1$, impossible. So every member has $n=1$:
$\Delta^-=e_a$ for one index $a$, and $\Delta^+=:m_a$ satisfies $|m_a|=k+1=N$ and $(m_a)_a=0$ ($\Delta^\pm$ have disjoint supports). We write
$\Delta_a=m_a-e_a$.

\emph{Step 2: value of a member.} $b\ge\Delta_a^-=e_a$ and $|b|=1$ force $b=e_a$, so $\mathcal D_{\Delta_a}=\{e_a\}$ with output $m_a$, and
$\Lambda(\Delta_a)=a_{e_a}\mult(m_a)$. Since $C$ is $S_d$-invariant, $Z_{e_a}$ and hence $a_1:=a_{e_a}$ do not depend on $a$.

\emph{Step 3: support disjointness.} For two members $\Delta_a,\Delta_{a'}$, Lemma~\ref{lem:S} gives
\[
\lVert\Delta_a-\Delta_{a'}\rVert_1=2(k+2)-2\,\ov(m_a,m_{a'})-2\,\ov(e_a,e_{a'}),
\]
so (ii), i.e.\ $\lVert\Delta_a-\Delta_{a'}\rVert_1\ge2(k+2)$, holds iff $\ov(m_a,m_{a'})=0$ and $\ov(e_a,e_{a'})=0$. Thus the indices $a$ of distinct members are
distinct, and the supports $S_a=\mathrm{supp}\,m_a$ are pairwise disjoint subsets of $[d]$. Put $s_a=|S_a|$. Then $1\le s_a\le s^\ast:=\min(N,d-1)$
($m_a\ne0$; $m_a$ has at most $N$ positive entries; $S_a\subset[d]\setminus\{a\}$), and $\sum_as_a\le d$.

\emph{Step 4: monotonicity of $f(s)=\Mb(N,s)/s$.} For integers $1\le s$ with $s+1\le\min(N,d)$, Lemma~\ref{lem:M}(M1) (applicable since $N>s$ and $s<d$) gives
\[
\frac{f(s+1)}{f(s)}=\frac{s}{s+1}\cdot\frac{\Mb(N,s+1)}{\Mb(N,s)}\ \ge\ \frac{2s}{s+1}\ \ge\ 1 .
\]
So $f$ is nondecreasing on $\{1,\dots,\min(N,d)\}$, and $s^\ast\le\min(N,d)$.

\emph{Step 5: numerator.} $\mult(m_a)\le\Mb(N,s_a)=s_af(s_a)\le s_af(s^\ast)$, so
\[\textstyle\sum_a\Lambda(\Delta_a)\le a_1f(s^\ast)\sum_as_a\le a_1\,d\,f(s^\ast).\]

\emph{Step 6: denominator.} Write $k=qd+s$ with $0\le s\le d-1$; $g\in C$ has $s$ coordinates $q+1$ and $d-s\ge1$ coordinates $q$, and $\mathcal D_g=\{e_j\}_{j\in[d]}$.
If $g_j=q$, then $g+e_j$ is balanced of size $N$, so $\mult(g+e_j)=\Mb(N,d)$. If $g_j=q+1$, moving one unit from coordinate $j$ (value $q+2$) to a
coordinate of value $q$ yields a balanced vector and multiplies $\mult$ by $(q+2)/(q+1)$ (move rule), so $\mult(g+e_j)=\Mb(N,d)\frac{q+1}{q+2}$. Hence
\[
\Lambda_0=a_1\sum_j\mult(g+e_j)=a_1\Mb(N,d)\Big[(d-s)+s\,\frac{q+1}{q+2}\Big],\qquad
\rho(F)\le\frac{d\,f(s^\ast)}{\Mb(N,d)\big[(d-s)+s\frac{q+1}{q+2}\big]} .
\]

\emph{Case A: $1\le k\le d-1$.} Then $q=0$, $s=k$, $N=k+1\le d$, so $\Mb(N,d)=N!$ and the bracket is $d-k/2$. Here $s^\ast=N$ if $k\le d-2$ and
$s^\ast=N-1$ if $k=d-1$; in both subcases $s^\ast\le N=\min(N,d)$, so $f(s^\ast)\le f(N)=N!/N=(N-1)!$ by Step 4. Hence
\[
\rho(F)\le\frac{d\,(N-1)!}{N!\,(d-k/2)}=\frac{d}{(k+1)(d-k/2)},
\]
and $(k+1)(d-\frac k2)-d=k\big(d-\frac{k+1}2\big)>0$ (as $k+1\le d$), so the bound is $<1$.

\emph{Case B: $k\ge d$.} Then $q\ge1$ and $N=k+1\ge d+1$, so $s^\ast=d-1$ and $f(s^\ast)=\Mb(N,d-1)/(d-1)$. The bracket is
$\ge d\frac{q+1}{q+2}$, and (M1) with $M=N>d-1=\ell$, $\ell<d$ gives $\Mb(N,d)\ge2\Mb(N,d-1)$. Hence
\[
\rho(F)\le\frac{d\,\Mb(N,d-1)/(d-1)}{2\Mb(N,d-1)\,d\,\frac{q+1}{q+2}}=\frac{q+2}{2(d-1)(q+1)}\le\frac{3}{4(d-1)}<1 .
\]
Cases A and B exhaust $k\ge1$ for every $d\ge2$ (for $d=2$: Case A is $k=1$, Case B is $k\ge2$).
\end{proof}
For $k=1$ the Case A bound is $d/(2d-1)$; at $d=3$ it equals $0.6$, the exact maximum found by the clique LP (research-phase computation
[Num]), so the bound is sharp there.

\subsection{Part (c): \texorpdfstring{$T=2$}{T=2}, all \texorpdfstring{$d\ge3$}{d>=3}, \texorpdfstring{$k\ge1$}{k>=1}}\label{app:T2}
Here $N=k+2$. The types of $|b|=2$ are A ($b=2e_i$, $c_b=1$, $d$ of them) and B ($b=e_i+e_j$, $c_b=2$, $\binom d2$ of them); by symmetry
$\kappa_b$ and $\theta_b$ depend only on the type, and in units of $M^\ast$, $W=d\sqrt{\kappa_A}+\binom d2\sqrt{2\kappa_B}$, $\theta_A=1/\sqrt{\kappa_A}$,
$\theta_B=\sqrt{2/\kappa_B}$.

\paragraph{T2.1 Structure.} By Lemma~\ref{lem:S}, $n\in\{1,2\}$ (S3), at most one member has $n=1$ (S2), and members with $n=2$ are pinned (S4).
The \emph{light} member (if any) is $\Delta_1=Q-e_i$ with $|Q|=k+1$, $Q_i=0$; its support is $\{2e_i\}\cup\{e_i+e_j\}_{j\ne i}$ with outputs
$Q+e_i$ and $Q+e_j$. A \emph{pinned} member is $\Delta=P-b$ with $|b|=2$, $|P|=N$, $\mathrm{supp}P\cap\mathrm{supp}b=\emptyset$, and value
$\theta_b\mult(P)$. The budgets of Lemma~\ref{lem:S} give: pinned--pinned ($n+n'-3=1$): $\ov(P,P')+\ov(b,b')\le1$, so the $b$'s are distinct
($\ov(b,b)=2$) and $\ov(P,P')\le1$; light--pinned ($n+n'-3=0$): $\ov(P,Q)=0$ and $b_i=0$.

\paragraph{T2.2 Light value.} With $t=|\mathrm{supp}Q|$ and the move rule $\mult(Q+e_x)=N\mult(Q)/(Q_x+1)$,
\begin{align*}
V_1(Q)&=\theta_A\mult(Q+e_i)+\theta_B\sum_{j\ne i}\mult(Q+e_j)\\&=N\mult(Q)\Big[\theta_A+\theta_B\Big((d-1-t)+\sum_{x\in\mathrm{supp}Q}\frac1{Q_x+1}\Big)\Big]\\
&\le\theta_AN\Mb(N-1,d-1)+(d-1)\theta_B\Mb(N,d-1).
\end{align*}

\paragraph{T2.3 Master bound.} Drop the light--pinned constraints. At most $d$ pinned members have type A, each of value $\le\theta_A\Mb(N,d-1)$.
For type B, $\ell=|\mathrm{supp}P|\le\min(d-2,N)$, so by Lemma~\ref{lem:J} ($r=2$),
$\mult(P)\le\Mb(N,\ell)\le\varphi_B\,s_2(P)$ with $\varphi_B=\max_{1\le\ell\le\min(d-2,N)}\Mb(N,\ell)/(\binom\ell2+[N>\ell])$, and
$\sum s_2(P)\le\binom{d+1}2$. Hence
\[
\rho(F)\le S(d,k):=\frac1W\Big[\theta_AN\Mb(N-1,d-1)+(d-1)\theta_B\Mb(N,d-1)+\theta_B\varphi_B\tbinom{d+1}2+d\,\theta_A\Mb(N,d-1)\Big].
\]
With $W\ge\binom d2\sqrt{2\kappa_B}$, $r_0=\Mb(N,d-1)/M^\ast$, $r_1=N\Mb(N-1,d-1)/M^\ast$, $r_\varphi=\varphi_B/M^\ast$:
\begin{equation}\label{eq:T2star}
S\le\frac2d\frac{r_0}{\kappa_B}+\frac{d+1}{d-1}\frac{r_\varphi}{\kappa_B}+\frac2{d-1}\frac{r_0}{\sqrt{2\kappa_A\kappa_B}}+\frac2{d(d-1)}\frac{r_1}{\sqrt{2\kappa_A\kappa_B}}.
\end{equation}

\paragraph{T2.4 Ratio lemmas.} Let $q_0=\lfloor k/d\rfloor$.
\begin{itemize}
\item[(K)] $\kappa_B\ge(\frac{q_0+1}{q_0+2})^2\ge\frac14$ and $\kappa_A\ge\frac{(q_0+1)^2}{(q_0+2)(q_0+3)}\ge\frac16$. \emph{Proof.} Fix $g\in C$ and let $y$ be
$g$ plus two units placed on smallest coordinates; $y$ is balanced of size $N$, so $\mult(y)=M^\ast$ and $\prod y_i!/\prod g_i!\ge(q_0+1)^2$.
Also $\prod(b+g)_i!/\prod g_i!\le(q_0+2)^2$ (type B) and $\le(q_0+2)(q_0+3)$ (type A). Since $\mult(b+g)/M^\ast=\frac{\prod y_i!/\prod g_i!}{\prod(b+g)_i!/\prod g_i!}$,
the bounds hold termwise in the average.
\item[(K0)] If $k\le d-2$: $M^\ast=N!$, $g$ is a $0/1$ vector with $k$ ones, $\kappa_B=\E\,2^{-X}\ge1-\E X/2=1-k/d$ with $X=g_i+g_j$ (as
$2^{-x}\ge1-x/2$ on $\{0,1,2\}$), $\kappa_A=\E[1/(2+g_i)!]=\frac12-\frac k{3d}$, and $r_\varphi\le1/\binom N2$. \emph{Proof of the last
bound.} If $N\le d-2$, $\ell=N$ gives $\Mb(N,N)/\binom N2=N!/\binom N2$. For $\ell\le N-1$, (M1) gives $\Mb(N,\ell)/(\binom\ell2+1)\le N!/g(\ell)$
with $g(\ell)=2^{N-\ell}(\binom\ell2+1)$; $g(\ell)/g(\ell+1)=2(\ell^2-\ell+2)/(\ell^2+\ell+2)\ge1\iff(\ell-1)(\ell-2)\ge0$, so $g$ is
nonincreasing, and $g(N-1)-\binom N2=\frac12(N^2-5N+8)>0$.
\item[(Moves)] By (M2) with $m=\lfloor N/d\rfloor$: $r_0\le R_m$ (one empty coordinate), $r_1\le mR_m$ ($r_1$ is $\mult(x)/M^\ast$ for an $x$ with a
coordinate equal to $1$), $r_\varphi\le R_m^2$. By (M1) when $N\ge d+1$: $r_0\le\frac12$ and
$r_\varphi\le\max_{\ell\le d-2}2^{\ell-d}/(\binom\ell2+1)=1/(4(\binom{d-2}2+1))$ (the maximand is $1/g$ with $N$ replaced by $d$).
\end{itemize}

\paragraph{T2.5 Lagrangian bound for the sharp cases.} Fix the shape of the light member $Q-e_i$, with $t=|\mathrm{supp}Q|$ and $o=d-1-t$.
Admissible pinned members have $b_i=0$ and $P\subset[d]\setminus\mathrm{supp}Q$, a ground set of size $d-t$, so Lemma~\ref{lem:J} gives
$\sum s_2(P)\le\binom{d-t+1}2=:\mathrm{cap}$. The possible $b$ fall into five classes:
\begin{center}\small
\begin{tabular}{lcc}\toprule
class of $b$ & count & coordinates available to $P$\\\midrule
$2e_j$, $j\in\mathrm{supp}Q$ & $t$ & $d-t$\\
$2e_j$, $j\notin\mathrm{supp}Q\cup\{i\}$ & $o$ & $d-t-1$\\
$e_j+e_l$, both in $\mathrm{supp}Q$ & $\binom t2$ & $d-t$\\
$e_j+e_l$, one in $\mathrm{supp}Q$ & $to$ & $d-t-1$\\
$e_j+e_l$, neither in $\mathrm{supp}Q$ & $\binom o2$ & $d-t-2$\\\bottomrule
\end{tabular}
\end{center}
Each $b$ carries at most one pinned member. For every $\lambda\ge0$ (weak duality),
\[
\sum_{\rm pinned}\theta_b\mult(P_b)\le\lambda\,\mathrm{cap}+\sum_b\max\Big(0,\ \max_P\big(\theta_b\mult(P)-\lambda s_2(P)\big)\Big),
\]
where the inner maximum runs over partitions $P$ of $N$ with at most the available number of parts ($s_2$ depends only on the shape). Thus
$\rho(F)\le\max(\rho_A,\rho_B)$, where $\rho_A$ is this Lagrangian divided by $W$ with no light member (all $d$ A-classes with $d-1$ available
coordinates and all $\binom d2$ B-classes with $d-2$, $\mathrm{cap}=\binom{d+1}2$), and $\rho_B=\max_Q[V_1(Q)+\text{Lagrangian}(Q)]/W$ over
partitions $Q$ of $k+1$ into at most $d-1$ parts.

\paragraph{T2.6 Region reduction.} Define (all from \eqref{eq:T2star})
\begin{align*}
E_1(d)&=\frac8d+\frac{d+1}{d-1}\,t_3(d)+\frac{2}{d-1}\Big(1+\frac1d\Big)\sqrt{12},\\ t_3(d)&=\max\Big\{\frac{1/3}{1-1/d},\frac{1/6}{1-2/d},\frac2{10},\frac{8}{(\lfloor d/2\rfloor+1)(\lfloor d/2\rfloor+2)}\Big\},\\
E_2^u(d)&=\frac8d+\frac{d+1}{d-1}\cdot\frac1{\binom{d-2}2+1}+\frac2{d-1}\Big(1+\frac1d\Big)\sqrt{12},\\
E_2(d,m)&=\frac{2}{d}\frac{\min(1,R_m)}{\kappa_B^{(m)}}+\frac{d+1}{d-1}\frac{\min\big(\frac1{4(\binom{d-2}2+1)},R_m^2\big)}{\kappa_B^{(m)}}\\
&\quad+\frac{\frac2{d-1}\min(1,R_m)+\frac2{d(d-1)}\min(1,mR_m)}{\sqrt{2\kappa_A^{(m)}\kappa_B^{(m)}}},
\end{align*}
with $\kappa_B^{(m)}=(\frac m{m+1})^2$ and $\kappa_A^{(m)}=\frac{m^2}{(m+1)(m+2)}$.
\begin{itemize}
\item \textbf{Region I} ($k\le d-2$, $d\ge40$). In \eqref{eq:T2star} use $r_0,r_1\le1$ ($M^\ast=N!$), $\kappa_B\ge\frac14$, $\kappa_A\ge\frac16$ and, for the
second term, $r_\varphi/\kappa_B\le1/(\binom N2\max(1-\frac kd,\frac14))\le t_3(d)$ (cases $k=1$, $k=2$, $3\le k\le d/2$, $k>d/2$). So $S\le E_1(d)$, and
every summand of $E_1$ is nonincreasing in $d$. $E_1(40)\le0.741501$ \checker{t2_interval/check_t2.py, step C1}.
\item \textbf{Region II} ($k\ge d-1$, $d\ge20$; so $N\ge d+1$). (K) with $q_0\ge0$, (Moves) $r_\varphi\le1/(4(\binom{d-2}2+1))$, and $r_0,r_1\le1$ give
$S\le E_2^u(d)$, nonincreasing in $d$. $E_2^u(20)\le0.790052$ \checker{t2_interval/check_t2.py, step C2}.
\item \textbf{Region II$'$} ($3\le d\le19$, $k\ge(m_0(d)+1)d-2$, so $m=\lfloor N/d\rfloor\ge m_0(d)+1$ and $N\ge d+1$). Since
$m=\lfloor(k+2)/d\rfloor\le q_0+1$, (K) with $q_0\ge m-1$ and (Moves) give $S\le E_2(d,m)$. $E_2$ is nonincreasing in $m$: the $\kappa^{(m)}$
increase, and $R_m$, $mR_m=m!/m^{m-1}$, $R_m^2$ decrease. With $m_0=4,3,3$ for $d=3,4,5$, $m_0=2$ for $6\le d\le16$ and $m_0=1$ for
$17\le d\le19$ (minimal with $E_2(d,m_0)<1$), $\max_dE_2(d,m_0)\le0.977398$ \checker{t2_interval/check_t2.py, step C3}.
\item \textbf{Finite remainder} ($d\le39$, $k\le d-2$; and $d\le19$, $d-1\le k\le(m_0+1)d-3$): 1021 pairs. $S(d,k)$ is evaluated in interval
arithmetic with $\mathrm{avg}_A,\mathrm{avg}_B$ exact rationals (closed form, checked against enumeration of $C$ for $d\le7$). $S<1$ except at 13 pairs;
the largest value below 1 is $\le0.996011$ at $(9,6)$ \checker{t2_interval/check_t2.py, step C4}.
\item \textbf{The 13 sharp pairs} $(d,k)\in\{(3,1)$, $(4,1)$, $(4,2)$, $(5,1)$, $(5,2)$, $(6,1)$, $(6,2)$, $(6,3)$, $(7,1)$, $(7,2)$, $(7,3)$, $(7,4)$, $(8,5)\}$: T2.5 evaluated in
interval arithmetic with stored multipliers $\lambda$ (any $\lambda\ge0$ is valid) gives $\rho\le0.679481$, attained at $(3,1)$
\checker{t2_interval/check_t2.py, step C5}.
\end{itemize}
By Lemma~\ref{lem:cover}(1) the regions cover all $d\ge3$, $k\ge1$, and the supremum of the bounds is $\le0.996011\le0.9961$. \qed

\subsection{Part (d): \texorpdfstring{$T=3$}{T=3}, all \texorpdfstring{$d\ge3$}{d>=3}, \texorpdfstring{$k\ge1$}{k>=1}}\label{app:T3}
Here $N=k+3$. The types of $|b|=3$ are A ($3e_i$, $c=1$, $d$ of them), B ($2e_i+e_j$, $c=3$, $d(d-1)$) and C ($e_i+e_j+e_l$, $c=6$,
$\binom d3$); in units of $M^\ast$, $\theta_X=\sqrt{c_X/\kappa_X}$ and $W=d\sqrt{\kappa_A}+d(d-1)\sqrt{3\kappa_B}+\binom d3\sqrt{6\kappa_C}$.

\paragraph{T3.1 Structure.} By Lemma~\ref{lem:S} (budget $n+n'-4$) the members are of three kinds: \textbf{L} ($n=1$): $Q-e_i$, $|Q|=k+1$, $Q_i=0$;
\textbf{M} ($n=2$): $P-c$, $|c|=2$, $|P|=k+2$, $\mathrm{supp}P\cap\mathrm{supp}c=\emptyset$; \textbf{Pin} ($n=3$): $R-b$, $|R|=N$,
$\mathrm{supp}R\cap\mathrm{supp}b=\emptyset$. L--L ($n+n'=2$) and L--M ($n+n'=3$) violate (S1), so a family has at most one L and never both an L and
an M. M--M has budget $0$: the $P$'s are pairwise disjoint. Pin--Pin has budget $2$: $\ov(R,R')\le2$ and the $b$'s are distinct ($\ov(b,b)=3$).
L--Pin and M--Pin constraints are dropped. Hence $\rho(F)\le(\max(V_L,V_M)+V_P)/W$, where $V_L$ bounds one L, $V_M$ all M's and $V_P$ all
pins.

\paragraph{T3.2 The three bounds} (units of $M^\ast$; $r_{A2}=\binom N2\Mb(N-2,d-1)/M^\ast$, $r_1=N\Mb(N-1,d-1)/M^\ast$, $r_0=\Mb(N,d-1)/M^\ast$).
\begin{itemize}
\item $V_L$: the support of $Q-e_i$ is $\{e_i+c':|c'|=2\}$ with outputs $Q+c'$: $b=3e_i$ gives $\mult(Q+2e_i)=\binom N2\mult(Q)$; $b=2e_i+e_j$ gives
$\mult(Q+e_i+e_j)=N\mult(Q+e_j)$; $b=e_i+2e_j$ and $b=e_i+e_j+e_l$ give vectors of size $N$ on $d-1$ coordinates. So
$V_L=\theta_Ar_{A2}+(d-1)\theta_B(r_1+r_0)+\binom{d-1}2\theta_Cr_0$.
\item $V_M$: the support of $P-c$ is $\{c+e_j\}_{j\in[d]}$, so its value is
\begin{align*}&\le\theta_{\max}\sum_j\mult(P+e_j)=\theta_{\max}N\mult(P)\Big[(d-\ell)+\sum_{x\in\mathrm{supp}P}\frac1{P_x+1}\Big]\\
&\le\theta_{\max}N\Mb(N-1,\ell)(d-\ell/2)\end{align*}
with $\ell=|\mathrm{supp}P|\le d-1$. The supports are disjoint, so $\sum\ell\le d$ and
$V_M=\theta_{\max}\,d\,N_f$ with $N_f=\max_{1\le\ell\le d-1}N\Mb(N-1,\ell)(2d-\ell)/(2\ell M^\ast)$.
\item $V_P$: at most $d$ A-pins, each $\le\theta_Ar_0$. For B- and C-pins, Lemma~\ref{lem:J} ($r=3$) gives $\sum s_3(R)\le\binom{d+2}3$ and
$\mult(R)\le\varphi_Xs_3(R)$ with $\varphi_X=\max_{\ell\le\min(d-a_X,N)}\Mb(N,\ell)/s_3^{\min}(\ell,N)$, $a_B=2$, $a_C=3$. So
$V_P=d\theta_Ar_0+\max(\theta_B\varphi_B,\theta_C\varphi_C)\binom{d+2}3/M^\ast$.
\end{itemize}
Define $S_3(d,k)=(\max(V_L,V_M)+V_P)/W$; then $\rho(F)\le S_3(d,k)$.

\paragraph{T3.3 Ratio lemmas.} Let $q_0=\lfloor k/d\rfloor$, $t=k/d$.
\begin{itemize}
\item[(K3)] $\kappa_A\ge\frac{(q_0+1)^3}{(q_0+2)(q_0+3)(q_0+4)}\ge\frac1{24}$, $\kappa_B\ge\frac{(q_0+1)^3}{(q_0+2)^2(q_0+3)}\ge\frac1{12}$,
$\kappa_C\ge(\frac{q_0+1}{q_0+2})^3\ge\frac18$: as (K), adding three units.
\item[(K3-0)] If $k\le d-3$ ($g$ is $0/1$, $M^\ast=N!$): $\kappa_C=\E2^{-X}\ge1-\frac32t$; $\kappa_B=\frac12\E[3^{-g_1}2^{-g_2}]\ge\frac12(1-\frac76t)$;
$\kappa_A=\E[1/(3+g_1)!]=\frac16(1-\frac34t)$.
\item[(J3)] If $N\ge4$, then $\Mb(N,\ell)/s_3^{\min}(\ell,N)\le N!/\binom N3$ for all $\ell\le N$. \emph{Proof.} Equality at $\ell=N$. At $\ell=N-1$:
$\Mb=N!/2$, $s_3^{\min}=\binom{N-1}3+N-2$, and $\binom N3\le2\binom{N-1}3+2(N-2)\iff\binom{N-1}2\le\binom{N-1}3+2(N-2)$, true for $N\ge6$
($\binom{N-1}3\ge\binom{N-1}2$) and for $N=4,5$ ($3\le5$, $6\le10$). For $\ell\le N-2$: $\Mb\le N!2^{\ell-N}$, $s_3^{\min}\ge\binom\ell3+\ell$, and
$h(\ell)=2^{N-\ell}(\binom\ell3+\ell)$ is nonincreasing ($h(\ell)/h(\ell+1)\ge1\iff\binom\ell3+\ell-1\ge\binom\ell2$, true for $\ell\ge1$);
finally $h(N-2)-\binom N3=4(\binom n3+n)-\binom{n+2}3=\frac n2(n^2-5n+10)\ge0$ with $n=N-2$.
\item[(J3$'$)] If $N\ge d+1$: $\varphi_X/M^\ast\le\max_{\ell\le d-2}2^{\ell-d}/(\binom\ell3+\ell)=1/(4(\binom{d-2}3+d-2))$ by (M1) and the same monotonicity.
\end{itemize}

\paragraph{T3.4 Region reduction.} Every step is either proved here or an evaluation by \nolinkurl{t3_master/t3_bounds.py}.
\begin{itemize}
\item \textbf{U1} ($d\ge100$, $1\le k\le d-3$, so $N\le d$, $M^\ast=N!$). Use $W\ge\binom d3\sqrt{6\kappa_C}$; $r_{A2},r_1,r_0\le1$; $N_f\le d/(N-1)$
(for $\ell\le N-1$, $N\Mb(N-1,\ell)/N!\le2^{\ell+1-N}$ and $\ell2^{N-1-\ell}\ge N-1$; for $\ell\ge N$ the ratio is $\le(2d-\ell)/(2\ell)\le d/N$);
$\varphi_X/M^\ast\le1/\binom N3$ (J3); and the larger of the (K3) and (K3-0) lower bounds for the $\kappa$'s, with $\theta_{\max}=\max(1/\sqrt{\kappa_A},\sqrt{3/\kappa_B},\sqrt{6/\kappa_C})$. This gives
\begin{align*}
S_3\le B_1(d,k)={}&\max\Big\{\frac{1}{\binom d3\sqrt{6\kappa_A\kappa_C}}+\frac{2(d-1)}{\binom d3\sqrt{2\kappa_B\kappa_C}}+\frac3d\frac1{\kappa_C},\ \frac{\theta_{\max}d^2}{(N-1)\binom d3\sqrt{6\kappa_C}}\Big\}\\
&+\frac{d}{\binom d3\sqrt{6\kappa_A\kappa_C}}+\max\Big\{\frac1{\sqrt{2\kappa_B\kappa_C}},\frac1{\kappa_C}\Big\}\frac{\binom{d+2}3}{\binom d3\binom N3}.
\end{align*}
For fixed $k$, $B_1$ is decreasing in $d$. For $k<12$ its value at $d=\max(100,k+3)=100$ is at most $0.301833$; for $k\ge12$, with the crude (K3)
constants, $B_1$ is also decreasing in $N$, so $B_1\le B_1^{\rm crude}(100,N=15)\le0.270687$ \checker{t3_master/t3_bounds.py}.
\item \textbf{U2} ($d\ge100$, $N\ge d+1$). (K3) crude constants, $r_1,r_{A2}\le1$, $r_0\le\frac12$ (M1), $N_f\le d/(d-1)$ (since
$N\Mb(N-1,\ell)\le\Mb(N,\ell+1)\le M^\ast2^{\ell+1-d}$ and $\ell2^{d-1-\ell}\ge d-1$) and (J3$'$) give a bound $B_2(d)$ of the same form, decreasing
in $d$, with $B_2(100)\le0.128161$ \checker{t3_master/t3_bounds.py}.
\item \textbf{F} ($3\le d\le99$, $N\ge\max(m_0(d)d,d+1)$, $m=\lfloor N/d\rfloor$). (K3) with $q_0\ge m-1$; (M2) with $t_0=0,1,2$:
$r_0\le\min(\frac12,R_m)$, $r_1\le\min(1,mR_m)$, $r_{A2}\le\min(1,m^2R_m/2)$ ($m\ge2$); $N\Mb(N-1,\ell)\le\Mb(N,\ell+1)\le M^\ast\min(2^{\ell+1-d},R_m^{d-1-\ell})$;
$\varphi_X/M^\ast\le\max_{\ell\le d-a_X}\min(2^{\ell-d},R_m^{d-\ell})/(\binom\ell3+\ell)$. This gives $S_3\le E_3(d,m)$, nonincreasing in $m$. The minimal
$m_0$ with $E_3(d,m_0)<1$ is $4$ ($d=3$), $3$ ($4\le d\le7$), $2$ ($8\le d\le14$), $1$ ($15\le d\le99$), and $\max_dE_3(d,m_0)\le0.968498$
\checker{t3_master/t3_bounds.py}.
\item \textbf{Finite} ($3\le d\le99$, $N<\max(m_0d,d+1)$): 4774 pairs, $S_3(d,k)$ evaluated with exact rational $\kappa$'s and interval square roots.
$S_3<1$ except at 18 pairs; the largest value below 1 is $\le0.970559$ at $(7,4)$ \checker{t3_master/t3_bounds.py}.
\item \textbf{The 18 residual pairs} $(d,k)\in\{(3,1)$, $(4,1)$, $(4,2)$, $(5,1)$, $(5,2)$, $(6,1)$, $(6,2)$, $(6,3)$, $(7,1)$, $(7,2)$, $(7,3)$, $(8,1)$, $(8,3)$, $(8,4)$, $(9,5)$, $(10,6)$, $(11,7)$, $(12,8)\}$:
each has an exact typed clique-LP certificate (Appendix~\ref{app:ca}) with bound $\le0.7717523470$ (attained at $(3,1)$) \checker{t3_clique_lp/check_torus_lp.py; cover: t3_master/check_t3_cover.py}.
\end{itemize}
\emph{Monotonicity used above.} In $B_1$ and $B_2$ every coefficient is a decreasing function of $d$: $\frac{2(d-1)}{\binom d3}=\frac{12}{d(d-2)}$,
$\frac d{\binom d3}$, $\frac{d^2}{\binom d3}=\frac{6d}{(d-1)(d-2)}$, $\frac{d^2}{(d-1)\binom d3}$, $\frac3d$ and $\binom{d+2}3/\binom d3=\frac{(d+2)(d+1)}{d(d-1)}$, and
$\binom{d+2}{3}/(\binom d3(\binom{d-2}3+d-2))$; the refined $\kappa$ bounds of (K3-0) increase as $t=k/d$ decreases, so $\theta_{\max}$ and every
$1/\sqrt{\kappa\kappa'}$ decrease; and $1/(N-1)$, $1/\binom N3$ decrease in $N$. In $E_3(d,m)$ the $\kappa^{(m)}$ increase with $m$ (so $W$ increases and
the $\theta$'s decrease), while $R_m$, $mR_m$ (ratio $(\frac m{m+1})^{m-1}\le1$), $m^2R_m/2$ for $m\ge2$ (ratio $(\frac m{m+1})^{m-2}\le1$; value $1$ at
$m=2$, matching the value $1$ used at $m=1$) and $\min(2^{\ell-d},R_m^{d-\ell})$ are nonincreasing; so $E_3$ is nonincreasing in $m$.

By Lemma~\ref{lem:cover}(2) the regions cover all $d\ge3$, $k\ge1$, and the supremum of the bounds is $\le0.970559\le0.9706$. \qed

\subsection{Coverage of the \texorpdfstring{$(d,k)$}{(d,k)} domain and exact comparison with \texorpdfstring{$1$}{1}}\label{app:coverage}
\begin{lemma}[Coverage]\label{lem:cover}
Let $m_0^{(2)}(d)$ and $m_0^{(3)}(d)$ be the tables of T2.6 and T3.4. (1) For $T=2$, every pair $(d,k)$ with $d\ge3$, $k\ge1$ lies in Region I,
Region II, Region II$'$ or the finite remainder of T2.6, and the hypotheses used for the bound of that region hold there. (2) For $T=3$, every pair
with $d\ge3$, $k\ge1$ lies in U1, U2, F or the finite set of T3.4, and the hypotheses of that region hold there.
\end{lemma}
\begin{proof}
(1) Fix $d\ge3$ and let $N=k+2$, $q_0=\lfloor k/d\rfloor$. If $d\ge40$: $k\le d-2$ is Region I and $k\ge d-1$ is Region II. If $20\le d\le39$: $k\le d-2$
is finite and $k\ge d-1$ is Region II. If $3\le d\le19$: $k\le d-2$ and $d-1\le k\le(m_0+1)d-3$ are finite, and $k\ge(m_0+1)d-2$ is Region II$'$.
In each case the listed ranges partition $\{k\ge1\}$, so there is no gap at any boundary. Hypotheses: Region I uses $k\le d-2$, i.e.\ $N\le d$ and
$M^\ast=N!$. Region II uses $N\ge d+1$, which is $k\ge d-1$. Region II$'$ uses $N\ge d+1$ and $m:=\lfloor N/d\rfloor\ge m_0+1$, both from
$N\ge(m_0+1)d\ge2d$; it also uses $q_0\ge m-1$, which holds since $N=q_0d+s+2$ with $0\le s\le d-1$ gives $N<(q_0+2)d$. Then (K) gives
$\kappa_B\ge(\frac{q_0+1}{q_0+2})^2\ge(\frac m{m+1})^2=\kappa_B^{(m)}$ and $\kappa_A\ge\kappa_A^{(m)}$ ($x\mapsto x/(x+1)$ and $x\mapsto x^2/((x+1)(x+2))$
are increasing), so $S\le E_2(d,m)\le E_2(d,m_0)$ by the monotonicity in $m$ proved in T2.6.

(2) Fix $d\ge3$, $N=k+3$. If $d\ge100$: $k\le d-3$ ($N\le d$) is U1 and $k\ge d-2$ ($N\ge d+1$) is U2. If $3\le d\le99$: $N<\max(m_0d,d+1)$ is the finite
set and $N\ge\max(m_0d,d+1)$ is F. Again these partition $\{k\ge1\}$. U1 uses $N\le d$ (so $M^\ast=N!$ and (K3-0) apply); for $k<12$ the bound
$B_1(d,k)$ is decreasing in $d$ (the refined $\kappa$ bounds increase as $t=k/d$ decreases, and every coefficient is decreasing in $d$), so it
suffices to evaluate it at $d=\max(100,k+3)=100$; for $k\ge12$ the crude-constant bound is decreasing in $d$ and in $N$, so it suffices to evaluate it
at $(d,N)=(100,15)$. U2 uses $N\ge d+1$ for (M1) and (J3$'$). F uses $N\ge d+1$, $m=\lfloor N/d\rfloor\ge m_0$, and $q_0\ge m-1$, which holds since
$N=q_0d+s+3<(q_0+2)d$ for $d\ge3$; then $S_3\le E_3(d,m)\le E_3(d,m_0)$.
\end{proof}

\paragraph{The checkers enumerate exactly these sets.} \nolinkurl{t2_interval/check_t2.py} evaluates $S(d,k)$ on
\texttt{[(d,k) for d in range(3,40) for k in range(1,d-1)]} $\cup$ \texttt{[(d,k) for d in range(3,20) for k in range(d-1,(m0[d]+1)*d-2)]},
i.e.\ on $\{3\le d\le39,\ k\le d-2\}\cup\{3\le d\le19,\ d-1\le k\le(m_0+1)d-3\}$, the finite remainder of Lemma~\ref{lem:cover}(1) (1021 pairs).
\nolinkurl{t3_master/t3_bounds.py} evaluates $S_3(d,k)$ on \texttt{range(1, max(m0[d]*d, d+1) - 3)} for $3\le d\le99$, i.e.\ $1\le k\le\max(m_0d,d+1)-4$,
which is $N<\max(m_0d,d+1)$ (4774 pairs). The tables $m_0^{(2)}$ ($4,3,3$ for $d=3,4,5$; $2$ for $6\le d\le16$; $1$ for $17\le d\le19$) and
$m_0^{(3)}$ ($4$ for $d=3$; $3$ for $4\le d\le7$; $2$ for $8\le d\le14$; $1$ for $15\le d\le99$) are recomputed by the checkers as the least $m$
with an interval upper bound $<1$ and compared with the stored data. An independent audit script, \nolinkurl{coverage_check.py} (shipped as a separate ancillary file next to \texttt{ca\_appendix.zip}, so the
archive hash is unchanged; SHA-256 \texttt{c42d11f652ed0bdd7daec55dd4bb9940319a8ff138c7f7d59956eacfa9ef0995}; it imports the shipped
checkers unmodified and takes the path of the unzipped archive as its argument; re-run independently in QH-Lab with output identical to the
original run), confirms on the box $3\le d<400$, $1\le k<1200$ that the enumerated sets equal the analytic definitions, that every pair lies
in some region, and that the region hypotheses ($N\ge d+1$, $N\le d$, $m\ge m_0$, $q_0\ge m-1$) hold pairwise; it also spot-checks the
monotonicity claims ($E_1,E_2^u$ in $d\le2000$; $E_2,E_3$ in $m\le60$, $m\le40$; $B_1,B_2$ in $d$) and evaluates the master bounds directly on
the boundary layers of every analytic region ($208$ pairs at $T=2$, max $S\le0.45272$; $206$ pairs at $T=3$, max $S_3\le0.33983$). These spot checks are
consistency checks [Num]; the monotonicity statements themselves are proved in T2.6, T3.4 and above.

\paragraph{Exact comparison with $1$; outward decimals.} Every certified inequality is decided exactly. In \nolinkurl{check_t2.py} and
\nolinkurl{t3_bounds.py} each bound is an \texttt{mpmath.iv} interval built from exact rationals (\texttt{fractions}), and the test is
\texttt{v.b < 1} on the upper endpoint, a binary floating-point number compared exactly with $1$. In \nolinkurl{check_torus_lp.py} the bound is a
\texttt{Fraction} (upper integer-square-root brackets in numerators, lower brackets in $\Lambda_0$), and the test \texttt{bound <= 1} is exact; the
bookkeeping script \nolinkurl{check_t3_cover.py} reads the printed LP bounds, which are rounded \emph{up} (\texttt{math.ceil} at $10^{-10}$). The decimals
printed by \nolinkurl{check_t2.py} and \nolinkurl{t3_bounds.py} are round-to-nearest; the constants quoted in this appendix are instead outward
(ceil) roundings of the interval upper endpoints: $0.741501$, $0.790052$, $0.977398$, $0.996011$, $0.679481$ ($T=2$); $0.301833$, $0.270687$, $0.128161$,
$0.968498$, $0.970559$, $0.7717523470$ ($T=3$). Their maxima give Theorem~\ref{thm:packing}(c) $\rho\le0.996011\le0.9961$ and (d) $\rho\le0.970559\le0.9706$.

\paragraph{What is computer-assisted.} In (c): the five evaluations C1--C5 (the 13 residual pairs at $T=2$ are closed by the interval
Lagrangian bound C5, not by clique-LP certificates). In (d): the constants of U1, U2 and F, the 4774 interval evaluations, and the 18 clique-LP
certificates. Everything else in this appendix, including (a), (b), the monotonicity statements and the
validity of every bound evaluated, is proved above.

\section{Proof of Lemma~\ref{lem:R}}\label{app:lemR}
\emph{Primal and dual.} With $J$ the Choi operator of a channel $\Sym^T\to\Sym^{T+k}$, $P_K(T,k)=\max\{\tr(JX):J\succeq0,\ \tr_{\rm out}J=I\}$,
where $X=\E_x v_xv_x^\dagger$ and $v_x=\bar\psi_x^{\otimes T}\otimes\psi_x^{\otimes(T+k)}$. The dual is $\min\{\tr Y:Y\otimes I\succeq X\}$.
$J=I/D_{T+k}$ is strictly feasible, so strong duality holds and the dual optimum is attained (Slater).

\emph{Sector form of $X$.} In the orthonormal occupation bases, $\psi_x^{\otimes m}=d^{-m/2}\sum_{|\gamma|=m}c_\gamma\,\omega^{x\,K\cdot\gamma}\ket\gamma$
with $c_\gamma=\sqrt{\mult(\gamma)}$ and $K\cdot\gamma=\sum_i\gamma_ik_i$. Hence $v_x=d^{-(2T+k)/2}\sum_{\alpha,\beta}c_\alpha c_\beta\,
\omega^{x\,K\cdot(\beta-\alpha)}\ket{\alpha,\beta}$, and $\E_x\omega^{x(u-u')}=[u\equiv u']$ gives
$X=\sum_{s\in\Z_q}\ket{w_s}\bra{w_s}$ with $w_s=d^{-(2T+k)/2}\sum_{K\cdot(\beta-\alpha)\equiv s}c_\alpha c_\beta\ket{\alpha,\beta}$.

\emph{Block structure of $Y$.} Let $V_x=\bar U_x^{\otimes T}\otimes U_x^{\otimes(T+k)}$ with $U_x=\mathrm{diag}(\omega^{xk_j})$. Then $V_xXV_x^\dagger=X$, so if $Y$
is feasible then so is $\bar U_x^{\otimes T}Y(\bar U_x^{\otimes T})^\dagger$, and so is their average $\bar Y$, with the same trace. $\bar Y$ commutes
with every $\bar U_x^{\otimes T}$, which acts on $\ket\alpha$ by the phase $\omega^{-x\,K\cdot\alpha}$; hence $\bar Y=\bigoplus_wY_w$ is block diagonal over the
classes $A_w=\{\alpha:K\cdot\alpha\equiv w\}$. We assume this form.

\emph{Reduction to scalar constraints.} The space splits into blocks $B_{w,\beta}=\mathrm{span}\{\ket{\alpha,\beta}:\alpha\in A_w\}$, on which
$Y\otimes I$ acts as $Y_w$. The vector $w_s$ meets block $B_{w,\beta}$ iff $K\cdot\beta-w\equiv s$, and there its component is
$d^{-(2T+k)/2}c_\beta v_w$ with $v_w=\sum_{\alpha\in A_w}c_\alpha\ket\alpha$; each block meets exactly one $w_s$. So $Y\otimes I-X$ is a direct sum, over
$s$, of $\bigoplus_{(w,\beta):K\cdot\beta-w\equiv s}Y_w$ minus the rank-one term $\ket{w_s}\bra{w_s}$. For $D\succeq0$ and a vector $u$,
$D\succeq uu^\dagger$ iff $u\in\mathrm{range}(D)$ and $u^\dagger D^+u\le1$ (Schur complement). Writing $g_w=v_w^\dagger Y_w^+v_w$ (and $g_w=+\infty$ if
$v_w\notin\mathrm{range}(Y_w)$), and using $\sum_{\beta:K\cdot\beta\equiv u}c_\beta^2=r_{T+k}(u)$ (each ordered $(T+k)$-tuple with sum $u$ is counted once
through its multiset), feasibility becomes
\[
d^{-(2T+k)}\sum_w r_{T+k}(s+w)\,g_w\le1\qquad\forall s\in\Z_q .
\]
Classes with $A_w=\emptyset$ do not occur, and for them $r_T(w)=0$.

\emph{Objective.} For $Y_w\succeq0$ and $v_w\in\mathrm{range}(Y_w)$, Cauchy--Schwarz gives
$\lVert v_w\rVert^4=(v_w^\dagger Y_w^{1/2}\,Y_w^{+1/2}v_w)^2\le(v_w^\dagger Y_wv_w)\,g_w\le\tr(Y_w)\lVert v_w\rVert^2g_w$, so
$\tr Y_w\ge\lVert v_w\rVert^2/g_w=r_T(w)/g_w$, with equality for $Y_w=v_wv_w^\dagger/g_w$. Since $g_w$ enters the constraints only through
$g_w$ itself, minimising $\tr Y$ is the same as minimising $\sum_wr_T(w)/z_w$ over $z_w=g_w>0$ subject to the displayed constraints. This is
Lemma~\ref{lem:R}.\qed

\section{Computer-assisted steps}\label{app:ca}

\paragraph{Package and integrity.}
All computer-assisted (CA) steps are packaged as \texttt{ca\_appendix.zip}, distributed as arXiv ancillary files and in the Zenodo record that accompanies this paper~\cite{QHLab}.
Its SHA-256 is
\begin{center}\footnotesize\texttt{e77e131bdd7becb0822ed582ccfc0a9e43b4b9183aeab37e65bec3e64d888b30}.\end{center}
The archive contains eight claim folders, a shared checker (\texttt{common/check\_torus\_lp.py}), the scripts that generated the certificates (\texttt{generators/}), and a \texttt{SHA256SUMS} file with 32 entries. The coverage audit \texttt{coverage\_check.py} (SHA-256 \texttt{c42d11f652ed0bdd7daec55dd4bb9940319a8ff138c7f7d59956eacfa9ef0995})
ships as a separate ancillary file next to the archive; see Appendix~\ref{app:coverage}.
After unzipping, \texttt{sha256sum -c SHA256SUMS} verifies every file.
Requirements: Python $\ge3.10$, \texttt{mpmath}, and \texttt{python-flint} 0.9.0. The clique-LP checker needs only the Python standard library.

\paragraph{Trust model.}
Every checker is standalone. It uses only exact rational arithmetic (\texttt{fractions}, with integer square roots for irrational quantities), interval arithmetic (\texttt{mpmath.iv}), or ball arithmetic (\texttt{python-flint} \texttt{arb}).
No checker calls an LP, SDP or ODE solver, NumPy or SciPy.
The certificates in \texttt{cert.json} and \texttt{*\_data.json} were produced by floating-point generators (HiGHS for the clique LPs) in \texttt{generators/}. They are \emph{data only}: weak duality makes any nonnegative multiplier valid, so correctness rests on the checkers alone.
The package was re-run independently on a separate machine as part of QH-Lab~\cite{QHLab} (v1.1.0, 49 automated tests), an independent Python re-implementation of the paper's tabulated values and
worked examples. The archive hash and all 32 internal hashes matched, and every checker's output was identical to the shipped \texttt{expected\_output.txt}. The three long clique-LP folders were run in resumable chunks, and their per-case lines were byte-identical.

\paragraph{The clique-LP checker (\texttt{check\_torus\_lp.py}).}
If two shifts $\Delta,\Delta'$ both lie within $\ell_1$-distance $T+k$ of a centre $c\in B_T$, then $\lVert\Delta-\Delta'\rVert_1\le2(T+k)$, which contradicts the separation condition (ii). Hence each such ball contains at most one family member, and for every $\mu\ge0$ on centres,
\[
\rho(F)=\sum_{\Delta\in F}r_\Delta\le\sum_c\mu_c+\sum_\Delta\max\Big(0,\,r_\Delta-\sum_{c:\lVert\Delta-c\rVert_1\le T+k}\mu_c\Big),
\qquad r_\Delta=\Lambda(\Delta)/\Lambda_0 .
\]
The certificate $\mu$ is a rational function of the centre's type.
The checker brackets the weights $a_b$ by integer square roots at scale $10^{30}$, using upper brackets in numerators and lower brackets in $\Lambda_0$.
It computes $Z_b$ exactly via elementary symmetric polynomials, enumerates shift types independently of the generator, and compares the final bound with $1$ exactly.
Members with $n(\Delta)=0$ cannot occur when $|F|\ge2$ (Lemma~S).
The checker also verifies $r_\Delta<1$ for unbalanced $\Delta\ge0$. This is redundant given~\cite[Prop.~X.5, Eq.~(148)]{Herasymenko2026}.

\subsection*{Folder-by-folder description}
Runtimes are wall-clock times on a single CPU core and are indicative only.

\paragraph{\texttt{t2\_interval} (Theorem T2).}
\emph{Claim:} steps C1--C5 of the proof of Theorem~T2 ($T=2$, all $d\ge3$, $k\ge1$; $\rho(F)\le0.9961$).
\emph{Method:} the claim splits into three parts.
\begin{itemize}
\item C1--C3 evaluate the closed-form region bounds $E_1(40)$, $E_2^u(20)$ and $E_2(d,m_0)$.
\item C4 evaluates $S(d,k)$ in interval arithmetic on 1021 finite pairs, with exact rational $\mathrm{avg}_A,\mathrm{avg}_B$. $S<1$ except at 13 pairs.
\item C5 evaluates the Lagrangian bound at the 13 residual pairs with stored multipliers $\lambda\ge0$.
\end{itemize}
\emph{Checker:} \texttt{python3 check\_t2.py} (reads \texttt{t2\_data.json}; uses Fractions and \texttt{mpmath.iv}).
\emph{Output:} C1 $0.741501$, C2 $0.790052$, C3 $\max 0.977398$, C4 $\max_{<1}0.996011$ at $(9,6)$, C5 $\max 0.679481$; final line \texttt{THEOREM T2 COMPUTATIONAL STEPS: PASS ; overall sup bound <= 0.9960} (this line is printed round-to-nearest; the
outward-rounded value is $0.996011$, and the PASS test itself compares the interval upper endpoint with $1$ exactly).
\emph{Runtime:} about 5\,s.

\paragraph{\texttt{t3\_master} (Theorem T3, analytic part).}
\emph{Claim:} the constants of the region reduction for Theorem~T3 ($T=3$, all $d\ge3$, $k\ge1$; $\rho(F)\le0.9706$). These are U1 (refined $0.301833$, crude $0.2707$), U2 ($0.1282$), and the $m_0$ table with $\max E_3=0.9685$.
\emph{Method:} the checker evaluates $S_3(d,k)$ on 4774 finite pairs with exact $\kappa$ and interval square roots. $S_3<1$ except at 18 pairs, and the maximum below $1$ is $0.97056$ at $(7,4)$. A bookkeeping script then checks that the 18 residual pairs are covered by \texttt{t3\_clique\_lp}.
\emph{Checker:} \texttt{python3 t3\_bounds.py 100 12 \&\& python3 check\_t3\_cover.py}. The cover check reads \texttt{t3\_bounds\_DU100.json}.
\emph{Output:} \texttt{t3\_bounds.py} ends with \texttt{analytic regions OK: True}; \texttt{check\_t3\_cover.py} prints \texttt{residual pairs: 18 ; max verified LP bound on them: 0.7717523470} and \texttt{THEOREM T3 (all d>=3, k>=1) COMPUTATIONAL STEPS: PASS}.
\emph{Runtime:} about 3\,s.

\paragraph{\texttt{t3\_clique\_lp} (Theorem T3, exact cases).}
\emph{Claim:} the packing inequality at $T=3$ for $d=2,\dots,12$ with $k\le20$ and for $d=13,\dots,20$ with $k\le12$ (316 cases). This range includes the 18 residual pairs of \texttt{t3\_master}.
\emph{Method:} typed clique-LP certificates with rational $\mu$.
\emph{Checker:} \texttt{python3 check\_torus\_lp.py cert.json}.
\emph{Output:} \texttt{CASES=316~~MAX\_BOUND<=0.7717523470~~ALL\_PASS=True}. The maximum is attained at $(d,k)=(3,1)$, where the exact value is $0.77175235$.
\emph{Runtime:} about 3\,min.

\paragraph{\texttt{clique\_lp\_d16\_Tk5} (certified region).}
\emph{Claim:} the packing inequality for $d\le16$ and $T+k\le5$ (150 cases).
\emph{Checker:} as above.
\emph{Output:} \texttt{CASES=150~~MAX\_BOUND<=0.8403003465~~ALL\_PASS=True}.
\emph{Runtime:} about 3\,s.

\paragraph{\texttt{d3\_T60} (certified region, $d=3$).}
\emph{Claim:} the packing inequality for $d=3$, $T\le40$ with $k\le3$, and $T=41,\dots,60$ with $k=1$ (140 cases).
\emph{Checker:} as above.
\emph{Output:} \texttt{CASES=140~~MAX\_BOUND<=0.9991923171~~ALL\_PASS=True}. The maximum is at $T=60$, consistent with $1-\max\rho\approx2.9/T^2$.
\emph{Runtime:} about 1\,min.

\paragraph{\texttt{t45\_clique\_lp} (evidence for Conjecture TT).}
\emph{Claim:} the packing inequality at $T=4$ for $d=3,\dots,10$ with $k\le12$ and $d=11,\dots,14$ with $k\le10$, and at $T=5$ for $d=3,\dots,9$ with $k\le6$ (178 cases).
\emph{Checker:} as above.
\emph{Output:} \texttt{CASES=178~~MAX\_BOUND<=0.8866406634~~ALL\_PASS=True}.
\emph{Runtime:} about 2\,min.

\paragraph{\texttt{su3\_cert} (SU(3) counterexample).}
\emph{Claim:} for $H=SU(3)$ acting on $V=\mathrm{Sym}^2\mathbb{C}^3$ and the Hessian group $G$ (216 elements modulo phases), at $(T,k)=(1,1)$ there is a $G$-orbit $X$ whose imposed moment is matched exactly to that of the $H$-orbit, yet $P_X-P_H\ge8.0399\times10^{-4}$, i.e.\ the finite orbit does not attain the continuous optimum.
\emph{Method:} ball arithmetic (\texttt{arb}, 256 bits). The checker certifies four things.
\begin{itemize}
\item Group closure by breadth-first search.
\item Commutant dimension and projector traces.
\item A Rayleigh lower bound on $P_X$ and a congruence-plus-Gershgorin upper bound on $P_H$.
\item An exactly moment-matched fiducial, found by Poincar\'e--Miranda on a preconditioned degree-$\le4$ map, with matching loss $\le8.15\times10^{-10}$.
\end{itemize}
\emph{Checker:} \texttt{python3 check\_su3.py} (reads \texttt{su3\_data.json}).
\emph{Output:} \texttt{gap\_lower [0.000803995920422 +/- 1.42e-16]~\textbar{}~matching\_loss\_upper [8.14587e-10 +/- 1.20e-17]~\textbar{}~PASS True}.
\emph{Runtime:} about 4\,s.

\paragraph{\texttt{spin1\_cert} (spin-1 counterexample).}
\emph{Claim:} the analogous strict gap for $SU(2)$ spin~1 with the octahedral group: an exactly moment-matched orbit with $P_X-P_H\ge6.437\times10^{-4}$.
\emph{Method:} ball arithmetic. The checker certifies group closure and the character check $\sum\chi^2/24=3$, the quadrature weights, and a Cholesky factorisation (minimum pivot $1.76\times10^{-12}$). It then certifies $P_X\ge0.516520734$ and $P_H\le0.515877020$, together with an intermediate-value sign change of the matching function, with loss bound $72\delta/\lVert f\rVert$.
\emph{Checker:} \texttt{python3 check\_spin1.py} (reads \texttt{spin1\_data.json}).
\emph{Output:} \texttt{gap\_lower [0.000643714380356 +/- 8.01e-17]~\textbar{}~matching\_loss\_upper [7.20000e-12 +/- 1e-22]~\textbar{}~PASS True}.
\emph{Runtime:} about 1\,s.

\subsection*{File hashes}
The six LP-folder copies of \texttt{check\_torus\_lp.py} are identical, including the copy in \texttt{common/}.
\begin{Verbatim}[fontsize=\scriptsize]
284885f682aa287943b2c356a01a61856ef8c19524c5491db3f65ca5530ef953  ./README.md
500fce4cabd10a639a457f56c28bb74d4ec3de8ae9c6b99f65b4d585c7a05da4  ./clique_lp_d16_Tk5/cert.json
abded76dbc9664403bef3f1d353e49dcf563c44d781bfcaee7c75ceab406b933  ./clique_lp_d16_Tk5/check_torus_lp.py
41de447e1afec21d4b7e453fe815417551b20043821da916d72afaaabd12b116  ./clique_lp_d16_Tk5/expected_output.txt
abded76dbc9664403bef3f1d353e49dcf563c44d781bfcaee7c75ceab406b933  ./common/check_torus_lp.py
4e1f176b9f97c453d60f773489a15b849edbd0c94d302b2bd68979da383527df  ./d3_T60/cert.json
abded76dbc9664403bef3f1d353e49dcf563c44d781bfcaee7c75ceab406b933  ./d3_T60/check_torus_lp.py
5fd48d2606e58d791a83493d76217d505f84d9a7c4f4d7eb0b1555ce16ee88d9  ./d3_T60/expected_output.txt
ee6920ce4bd8a0cd6edd90e544fce29ec03cfbbb9cba01cce9e47130afcf1197  ./generators/export_spin1.py
477f525cf30c19038bdbaf8cc8ff461a2dd512004a191b94b18699169f18ddb0  ./generators/export_su3.py
a48637ace29fbc91df98df07c256343b52d97923713159fbb62ca2e4ba825d27  ./generators/gen_torus_cert.py
b26f05d6859912e7be78c57b585e8a3a759bde9cceb864de784db722ef229d22  ./spin1_cert/check_spin1.py
98b6358ed32b875426b787b4d93b8c56bc248ec77006fd79a74411c6af232dc9  ./spin1_cert/expected_output.txt
496edcc41b5d150ee90c63fcc29b31b02be017b7d11143959f5749a0ca0a211c  ./spin1_cert/spin1_data.json
5cbce2c9ec3e80c90c817b18edcbfc6c49e7d76d40adb2dcba7de4585afaf75b  ./su3_cert/check_su3.py
2d3c3b46dc642108bbb9c5aba17085c691f83749b246e0ddc247b813dacd8ae5  ./su3_cert/expected_output.txt
54c1af221db10d21f91acd5f29a30ad47502fb7547e4a8c668b171c2ea1b0741  ./su3_cert/su3_data.json
7360fcb026a4cda4d0d39f2c7c1083518b8248c1c6caad9b0118b4436629e89e  ./t2_interval/check_t2.py
90ab1ee6f5472d41083fb15fce1771f9274145dbb72db688cd8735c75d40ccf4  ./t2_interval/expected_output.txt
4378b2d4ab43a82bade4c561fc1afc2723660e96ed0dc37ae6540b9dd9015481  ./t2_interval/make_t2_data.py
0912d57c75d71e7d445dea4670ef0dcfda87733cfd7d7c70fe55e4ae418045e0  ./t2_interval/t2_data.json
67dc7fbdf256efa51084eae5f6f3a742a625b2306783727b498701e360e27b0b  ./t3_clique_lp/cert.json
abded76dbc9664403bef3f1d353e49dcf563c44d781bfcaee7c75ceab406b933  ./t3_clique_lp/check_torus_lp.py
baefcd6d59d5e11c9d0551a8f6fe19e96bcb38bd06d7ac3dcbee89d560112272  ./t3_clique_lp/expected_output.txt
85399ef85a3e8cfeac493191cfedbd367fb5a70dd0641c3b0108bb2feab88566  ./t3_master/check_t3_cover.py
6d701bf17296510b8b30e84a34a0e12e99bebc307f4aa3f6515bd1ce59dc6859  ./t3_master/expected_cover_output.txt
00b9ff2e859e361127357a94639a7a1f038eea6919d69ace09613034976bb9d5  ./t3_master/expected_output.txt
3dcbb1beee39b34f72436241aacf60c426fe231dfb17b0caa74db7c3e6335235  ./t3_master/t3_bounds.py
59338e4eea8d368305114056617ff9ca07de8f006d1a2503cbb38f5c00c60b96  ./t3_master/t3_bounds_DU100.json
077f931afb2e786692ab2733cff53c556f4615e5e33c2d16c998973d446b7d95  ./t45_clique_lp/cert.json
abded76dbc9664403bef3f1d353e49dcf563c44d781bfcaee7c75ceab406b933  ./t45_clique_lp/check_torus_lp.py
6278e682a9d6a6e9d2c71121d242fa7aead8ee6e1c98f964ec32e1a5281a3506  ./t45_clique_lp/expected_output.txt
\end{Verbatim}

\end{document}
